\documentclass[11pt,a4paper]{article}

\usepackage[utf8]{inputenc}
\usepackage[T1]{fontenc}
\usepackage{amsmath,amssymb}
\usepackage{graphicx}
\usepackage{booktabs}
\usepackage{hyperref}
\usepackage[margin=1in]{geometry}
\usepackage{caption}
\usepackage{subcaption}
\usepackage{float}
\usepackage{xcolor}
\usepackage{xspace}

%%% Custom commands (harmonized across post_preliminary_updates.tex
%%% and truth_vertex_reco_check.tex)
\newcommand{\ET}{\ensuremath{E_{\text{T}}}\xspace}
\newcommand{\eT}{\ensuremath{E_{\text{T}}}\xspace}
\newcommand{\pT}{\ensuremath{p_{\text{T}}}\xspace}
\newcommand{\etg}{\ensuremath{E_{\mathrm{T}}^{\gamma}}\xspace}
\newcommand{\isoET}{\mbox{$E_{\text{T}}^{\text{iso}}$}\xspace}
\newcommand{\isoETtruth}{\ensuremath{E_\text{T}^\text{iso,truth}}\xspace}
\newcommand{\wetacog}{\ensuremath{w_{\eta}^{\mathrm{CoG}}}\xspace}
\newcommand{\chidf}{\ensuremath{\chi^2/\mathrm{ndf}}\xspace}
\newcommand{\DR}{\ensuremath{\Delta R}\xspace}
\newcommand{\zvtx}{\ensuremath{z_{\mathrm{vtx}}}\xspace}
\newcommand{\pp}{\mbox{$p$+$p$}\xspace}
\newcommand{\GeV}{\ifmmode {\mathrm{\ Ge\kern -0.1em V}}\else\textrm{Ge\kern -0.1em V}\fi\xspace}
\newcommand{\nown}{\ensuremath{n_{\mathrm{owned}}}\xspace}
\newcommand{\dietag}{\ensuremath{\Delta i_\eta}\xspace}
\newcommand{\diphig}{\ensuremath{\Delta i_\phi}\xspace}
\newcommand{\splitnode}{\texttt{CLUSTERINFO\_CEMC}\xspace}
\newcommand{\nosplitnode}{\texttt{CLUSTERINFO\_CEMC\_NO\_SPLIT}\xspace}

\title{\LARGE\textbf{Double-Interaction (Pileup) Effects in PPG12:\\
Efficiency, Shower Shapes, Truth-Vertex Reweighting,\\
Cluster Size, and Energy Response}\\[6pt]
\large Consolidated Note ---
\textbf{\textit{sPHENIX} Internal} ---
\texttt{sPH-ppg-2024-012}}
\author{PPG12 Analysis Team}
\date{April 18, 2026}

\begin{document}

\maketitle

\begin{abstract}\noindent
Run-24 \pp data at $\sqrt{s} = 200$\,\GeV contain two crossing-angle periods
(0\,mrad and 1.5\,mrad) with very different pileup conditions, where two
independent collisions in the same bunch crossing shift the reconstructed
vertex and distort cluster kinematics, shower shapes, and ABCD templates.
This note consolidates the PPG12 treatment of double interactions (DI) across
four lines of evidence: (i)~luminosity-weighted MBD-Poisson pileup fractions
($22.4\%$ cluster-weighted at 0\,mrad, $7.9\%$ at 1.5\,mrad);
(ii)~the DI efficiency study, which decomposes the apparent $94\%\to 31\%$
double-MC drop into an $88\%$ matching artifact (removed by switching from
$\DR$-based to track-ID-only matching) and $12\%$ genuine cluster loss,
leaving a $9.4\%$ relative efficiency reduction (52.0\%/57.4\%) on the
mixed-0\,mrad integrated photon-ID chain and a 2.8\% integrated
cross-section impact at 1.5\,mrad (reco 1.6\% + BDT 1.0\% + iso 0.2\%);
(iii)~an iterative data-driven truth-vertex reweight $w(z)$ factorized as
$w(z_h)\,w(z_\mathrm{mb})$, which closes the individual DI truth-vertex
distributions that the production $f(z_r)$ structurally cannot;
(iv)~data-vs-mixed-MC shower-shape validation with up to $20.8\times$
\chidf improvements at 0\,mrad for vertex-sensitive variables.
The nominal recipe is: single-MC efficiency with track-ID matching, the
iterative truth-vertex reweight on both periods, mixed MC for shower-shape
validation only. The $+1.4\%$ 1.5\,mrad cross-section bias under
$f(z_r)$ alone is removed inside the $|z_h|<60$\,cm fiducial, with per-bin
closure at $1$--$2\%$.
\end{abstract}

\tableofcontents
\clearpage

%% ==================================================================
\section{Introduction}
\label{sec:intro}

In \pp collisions at RHIC, two separate collisions can occur in the same
bunch crossing (pileup, also referred to here as double interaction, DI).
For PPG12's isolated prompt-photon cross-section measurement at
$\sqrt{s}=200$\,\GeV using sPHENIX Run~24 data, pileup is not just a
rate correction: it shifts the reconstructed MBD vertex to the average of
the two collision points, which in turn distorts cluster pseudorapidity
via the EMCal-radius geometric mapping, biases the truth-reco matching
that underlies the efficiency calculation, broadens \wetacog at the
shower-shape level, enters the ABCD background templates through
DI-broadened jet clusters, and --- if uncorrected --- biases the truth
fiducial acceptance used for the efficiency denominator at the
percent level.

The analysis uses two crossing-angle periods (split at run 51274):
\begin{itemize}
\item \textbf{0\,mrad}, runs 47289--51274, $\mathcal{L}\simeq 32.7$\,pb$^{-1}$,
      broad luminous region ($\sigma_z\!\simeq\!54$\,cm), high pileup
      ($\sim\!22\%$ cluster-weighted DI);
\item \textbf{1.5\,mrad} (nominal published period),
      runs 51274--54000, $\mathcal{L}\simeq 16.9$\,pb$^{-1}$,
      Piwinski-narrowed luminous region ($\sigma_z\!\simeq\!22$\,cm),
      modest pileup ($\sim\!8\%$ cluster-weighted DI).
\end{itemize}

The PPG12 DI programme consists of two complementary work streams:
\begin{enumerate}
\item \textbf{Full-GEANT DI MC blending}: dedicated
  \texttt{photon10\_double} and \texttt{jet12\_double} samples, produced
  via full detector simulation of two overlapping \pp collisions, are
  mixed with the nominal single-interaction samples at the physics
  fractions of Table~\ref{tab:pileup_fractions}. This blend provides both
  the DI efficiency study (\S\ref{sec:eff_impact}) and the
  shower-shape data--MC validation
  (\S\ref{sec:shower_shape}).
\item \textbf{Truth-vertex reweighting}: the production reco-level weight
  $f(z_r)=D(z_r)/M(z_r)$ is structurally insufficient to close the
  individual truth vertices of double-interaction events; an iterative
  data-driven reweight $w(z)$ applied per truth vertex
  ($w(z_h)$ for single MC, $w(z_h)\,w(z_\mathrm{mb})$ for doubles) is
  introduced in \S\ref{sec:truth_vertex} and removes the $+1.4\%$
  cross-section bias at 1.5\,mrad.
\end{enumerate}

A third, toy-level vertex-shift simulation
(\texttt{DoubleInteractionCheck.C}) exists for cross-checks on single MC
and is described in \texttt{.claude/rules/double-interaction.md}; this note
focuses exclusively on the full-GEANT and truth-vertex-reweight work
streams that feed the nominal analysis.

The remainder of the note covers: pileup fractions (\S\ref{sec:fractions}),
DI MC production (\S\ref{sec:production}), the vertex-shift mechanism
(\S\ref{sec:vtx_shift}), the efficiency impact with the full DI
decomposition (\S\ref{sec:eff_impact}), shower-shape data--MC validation
(\S\ref{sec:shower_shape}), truth-vertex reweighting
(\S\ref{sec:truth_vertex}), cluster size under DI (\S\ref{sec:cluster_size}),
energy response (\S\ref{sec:energy_response}),
eta-edge migration (\S\ref{sec:eta_edge}), and the summary
(\S\ref{sec:summary}). Appendices cover the DI reco--truth kernel,
discarded reweighting methods, the baseline TruthVertexRecoCheck study,
and a file/code inventory.

%% ==================================================================
\section{Pileup Fractions}
\label{sec:fractions}

The mean number of interactions per bunch crossing, $\mu$, is derived
from the Poisson-corrected MBD coincidence rate:
\begin{equation}
    \mu = \frac{\sum_{\text{runs}} N_{\text{MBD,corr}}}{r_B \cdot \sum_{\text{runs}} \Delta t},
\end{equation}
where $r_B = 111 \times 78{,}200 = 8.68 \times 10^{6}$\,Hz is the RHIC
bunch-crossing rate. The event-level interaction fractions under the
zero-truncated Poisson are
\begin{equation}
    f_1 = \frac{\mu\, e^{-\mu}}{1 - e^{-\mu}},\qquad
    f_2 = \frac{(\mu^2/2)\, e^{-\mu}}{1 - e^{-\mu}},\qquad
    f_{\geq 3} = 1 - f_1 - f_2.
\end{equation}
Fractions are computed per run and {luminosity}-weighted (not
simple-averaged) to avoid Jensen's inequality bias. Cluster-weighted
fractions account for the higher cluster multiplicity of
multi-interaction events:
\begin{equation}
    w_{\text{single}} = \frac{f_1}{\sum_k k\, f_k},\qquad
    w_{\text{double}} = \frac{2\, f_2 + 3\, f_{\geq 3}}{\sum_k k\, f_k},
\end{equation}
with triple-and-higher interactions folded into the double bin for the
purposes of the GEANT blend (the $\tt photon10\_double$ /
$\tt jet12\_double$ productions model the two-collision case exactly;
extra collisions are rare enough that this approximation is well within
the bin-by-bin closure tolerance).

\begin{table}[H]
    \centering
    \caption{Luminosity-weighted pileup fractions and cluster-level
      mixing weights from \texttt{calc\_pileup\_range.C}. Triple+ folded
      into $f_2$ for blending. $\mu_{\text{corr}}$ is the MBD-corrected
      mean interactions per bunch crossing.}
    \label{tab:pileup_fractions}
    \begin{tabular}{lcccccc}
        \toprule
        Crossing angle & Run range & $\mu_{\text{corr}}$ & $f_1$ & $f_2$ & $w_{\text{single}}$ & $w_{\text{double}}$ \\
        \midrule
        0~mrad   & 47289--51274 & 0.24 & 0.879 & 0.111 & 0.776 & \textbf{0.224} \\
        1.5~mrad & 51274--54000 & 0.08 & 0.960 & 0.039 & 0.921 & \textbf{0.079} \\
        \bottomrule
    \end{tabular}
\end{table}

The cluster-weighted fractions $w_{\text{double}}$ (last column of
Table~\ref{tab:pileup_fractions}) are the numbers used throughout
this note: $0.224$ at 0\,mrad and $0.079$ at 1.5\,mrad. They
multiply the per-event cross-section weight of each sample in the
full-GEANT blending pipeline (\S\ref{subsec:two_pass}) and set the
double-interaction mixing fraction $f_d$ in the iterative truth-vertex
forward model (\S\ref{subsec:method}).


%% ==================================================================
\section{DI MC Production}
\label{sec:production}

The full-GEANT DI samples used throughout this note are
\texttt{photon10\_double} (photon signal, truth $p_T^\gamma \in
[10,100]$\,\GeV with window $[14,30]$\,\GeV dominant) and
\texttt{jet12\_double} (QCD jet, $p_T^{\text{jet}}\in[14,21]$\,\GeV),
produced in the run28 production area
\texttt{anatreemaker/macro\_maketree/sim/run28/}. The double-interaction
events are generated via the Fun4All embedding flag~2
(pythia-in-hijing), carrying a primary hard-scatter vertex ($z_h$) and
a secondary minimum-bias vertex ($z_\mathrm{mb}$). Both vertex values
are available in the output slimtree through the branches
\texttt{vertexz\_truth} ($z_h$) and \texttt{vertexz\_truth\_second}
($z_\mathrm{mb}$, read in C++ as \texttt{vertexz\_truth\_mb}).

\paragraph{Known production caveats.} Two independent pipeline reviews
(\texttt{condor\_pipeline\_review\_run28.md}, findings \#7 and \#19;
the BDT training audit of Apr 2026) flagged the following issues that
the reader should be aware of:

\begin{itemize}
\item \textbf{Hardcoded \texttt{total\_line=10000} on the pre-submit
  stage}: the per-sample \texttt{run\_condor.sh} templates in 12 of 14
  sample directories hardcode a 10,000-line filelist cut-off, while
  \texttt{photon10\_double} (9998 lines) and \texttt{jet12\_double}
  (9997) sit just below the cutoff. The final 2--3 jobs per double
  sample therefore read empty filelists and produce empty
  \texttt{caloana.root} files that are silently merged into the combined
  output by \texttt{hadd\_combined.sh}. Two jet samples (\texttt{jet5},
  \texttt{jet8}) already patched this with a dynamic \texttt{wc -l}
  look-up; the remaining samples have not.
\item \textbf{Missing \texttt{vertexz\_truth\_mb} branch}: the
  \texttt{photon10\_double} and \texttt{jet12\_double}
  \texttt{bdt\_split.root} files currently downstream of
  \texttt{apply\_BDT.C} were produced before the
  \texttt{vertexz\_truth\_mb} branch was added to the anatreemaker output.
  Consumers downstream (notably \texttt{ShowerShapeCheck.C} and
  \texttt{RecoEffCalculator\_TTreeReader.C}, which both carry an
  optional-pointer guard) silently drop the secondary-vertex factor
  $w(z_\mathrm{mb})$ for those two samples. The impact on cluster-size
  means and shower-shape moments is expected to be below the percent
  level at 1.5\,mrad (where the DI blend fraction is only 7.9\%) and
  sub-percent at 0\,mrad (where the MC beam profile is already close to
  the luminous region). Regenerating these two
  \texttt{bdt\_split.root} files from the reprocessed
  \texttt{combined.root} (which does carry \texttt{vertexz\_truth\_mb})
  removes this caveat; it is a follow-up.
\item \textbf{Embedding flag buried in 3 of 14 Fun4All macro copies}:
  \texttt{TruthJetInput::add\_embedding\_flag(2)} is set in the Fun4All
  macros of \texttt{photon10\_double}, \texttt{jet12\_double}, and
  \texttt{jet8} (for distinct physics reasons); any consolidation of the
  14-way macro duplication should preserve this flag explicitly on the
  DI productions.
\end{itemize}

None of these issues blocks the DI programme at its current precision; the
cross-section bias from the missing $w(z_\mathrm{mb})$ factor is
sub-percent and well below the existing systematic budget at 1.5\,mrad.
The pipeline consolidation is tracked separately in
\texttt{condor\_pipeline\_review\_run28.md}.


%% ==================================================================
\section{Vertex-Shift Mechanism}
\label{sec:vtx_shift}

In a double-interaction event the MBD reconstructs a single vertex
whose best estimator is the average of the two collision points:
\begin{equation}
    z_{\text{vtx}}^{\text{reco}} \;\approx\; \frac{z_{\text{vtx},1} + z_{\text{vtx},2}}{2}.
\end{equation}
The RHIC luminous region has $\sigma_\mathrm{lum} \approx 22$\,cm at
1.5\,mrad and $\sigma_\mathrm{lum} \approx 54$\,cm at 0\,mrad, while
the MC generator was produced with a wider
$\sigma_\mathrm{gen} \approx 65$\,cm common to both periods. The
two vertices in a DI event are drawn from the generator distribution
(MC) or from the luminous region (data), so the displacement between
the primary truth vertex and the reconstructed vertex,
$\Delta z \equiv z_{\text{vtx}}^{\text{reco}} - z_{\text{vtx}}^{\text{truth}}$,
can reach several tens of cm in either case. The cluster pseudorapidity shift induced at
the EMCal-barrel radius $R_{\text{EMCal}} = 93.5$\,cm follows directly
from the cluster geometry:
\begin{equation}
    \Delta\eta \;\approx\; -\frac{\Delta z}{R_{\text{EMCal}} \cosh\eta}.
    \label{eq:deta_vs_dz}
\end{equation}
At $\eta = 0$ this reduces to $|\Delta\eta| \approx |\Delta z|/93.5$\,cm,
so a matching cut of $\DR<0.1$ fails whenever
$|\Delta z| \gtrsim 9.35$\,cm at midrapidity. This is the critical
observation that drives the efficiency interpretation in
\S\ref{sec:eff_impact}: pileup events populate $|\Delta z|$ far beyond
this threshold.

\begin{table}[H]
    \centering
    \caption{Vertex-shift distributions in single- and double-interaction
      photon10 MC. The critical threshold $|\Delta z| = 9.35$\,cm maps
      onto $|\Delta\eta| = 0.1$ at $\eta = 0$ via
      Eq.~\eqref{eq:deta_vs_dz}.}
    \label{tab:vertex_shift}
    \begin{tabular}{lcc}
        \toprule
        Quantity & Single MC & Double MC \\
        \midrule
        RMS of $\Delta z$ [cm] & 7.6 & 34.5 \\
        Fraction with $|\Delta z| > 9.35$\,cm & --- & 83\% \\
        \bottomrule
    \end{tabular}
\end{table}

The matching-failure rate inherits the same deterministic geometric
origin, saturating to $100\%$ for vertex displacements well past the
$9.35$\,cm threshold:

\begin{table}[H]
    \centering
    \caption{$\DR<0.1$ matching-failure rate vs.\ $|\Delta z|$ in
      double-interaction MC.}
    \label{tab:failure_vs_dz}
    \begin{tabular}{lc}
        \toprule
        $|\Delta z|$ range [cm] & Matching-failure rate \\
        \midrule
        $< 10$ & 9\% \\
        $10$--$15$ & 77\% \\
        $15$--$20$ & $\sim$100\% \\
        $> 20$ & 100\% \\
        \bottomrule
    \end{tabular}
\end{table}

Figure~\ref{fig:fail_vs_dz} shows the same quantity as a continuous
curve in single-interaction MC, where the rare tail of large $|\Delta z|$
events gives a clean step at the geometric threshold.

\begin{figure}[H]
    \centering
    \includegraphics[width=0.75\textwidth]{../plotting/figures/deltaR_study/fail_fraction_vs_dz.pdf}
    \caption{$\DR<0.1$ failure fraction as a function of $|\Delta z|$ in
      single-interaction MC. The sharp step at ${\sim}9.35$\,cm matches
      the prediction from Eq.~\eqref{eq:deta_vs_dz}, confirming vertex
      mismatch as the sole driver of $\DR$ matching failure; cluster
      quality is essentially unchanged across the step.}
    \label{fig:fail_vs_dz}
\end{figure}


%% ==================================================================
\section{Efficiency Impact}
\label{sec:eff_impact}

\subsection{Apparent Efficiency Drop Under \DR Matching}
\label{subsec:apparent_drop}

When the \texttt{photon10\_double} sample was first passed through the
legacy efficiency pipeline (which required a $\DR<0.1$ geometric match
on top of the GEANT track-ID to associate a reconstructed cluster with
a truth photon), the reconstruction efficiency dropped dramatically:
from $\sim\!94\%$ in single-interaction MC to $\sim\!31\%$ in the
double-interaction sample. Taken at face value this would imply that
pileup destroys two-thirds of the reconstructed cluster sample --- a
result in obvious tension with the $\sim\!1$\% fiducial-acceptance bias
expected from the pileup fraction alone.

A dedicated investigation showed that the apparent drop was
overwhelmingly a \emph{truth--reco matching artifact} caused by the
vertex-displacement mechanism of \S\ref{sec:vtx_shift}. The single
decisive observation is that the matching failure rate scales as
expected from Eq.~\eqref{eq:deta_vs_dz}: in double MC, 83\% of events
have $|\Delta z|>9.35$\,cm, of which 77\%+ fail the $\DR<0.1$ cut.
The cluster is physically present at the correct track-ID; the
reconstructed $\eta$ is merely offset from the true $\eta$ by the
vertex shift. Switching the matching strategy from $\DR$-geometric to
track-ID-only recovers those clusters and eliminates the artifact.
This section decomposes the apparent drop, justifies the removal of
$\DR<0.1$ from the production pipeline, and reports the factorized
per-stage efficiency under the corrected matching.

\subsection{Decomposition: Matching Artifact vs.\ Genuine Cluster Loss}
\label{subsec:decomposition}

For truth photons that failed the $\DR<0.1$ match in the double MC,
the failure mode was classified by searching for a cluster with the
correct \texttt{trkID} irrespective of $\DR$:
\begin{itemize}
\item \textbf{\DR failure (88\%):} a cluster with the correct
  \texttt{trkID} exists, but sits at $\DR>0.1$ due to vertex
  displacement. The signal-photon cluster is fully reconstructed and
  carries the right energy; it simply appears at the wrong
  reconstructed $\eta$ when the cluster position is taken with respect
  to the displaced MBD vertex.
\item \textbf{Missing cluster (12\%):} no cluster carries the matching
  \texttt{trkID}. This is genuine cluster loss --- the photon shower
  merged into a neighbouring jet shower or was absorbed by the
  higher-multiplicity environment.
\item \textbf{Extra truth photons ($\sim 0\%$):} the second collision
  does not produce additional truth-level prompt photons at the level
  relevant for PPG12. A per-event check on 100k photon10 events shows
  $0.493$ truth photons/event in single MC and $0.495$ in double MC
  (statistically identical); the $\geq 2$-denominator fraction is
  $0.22\%$ vs.\ $0.21\%$; see Table~\ref{tab:truth_content}.
\end{itemize}

\begin{table}[H]
    \centering
    \caption{Truth-level prompt photon content in photon10 samples
      (100k events each). Denominator: $|\eta^\gamma|<0.7$,
      $\isoETtruth<4$\,\GeV. The slight increase in mean truth
      $\isoETtruth$ in double MC (0.26 vs.\ 0.17\,\GeV) reflects
      additional hadrons from the second collision falling within the
      isolation cone, but remains well below the 4\,\GeV fiducial
      threshold.}
    \label{tab:truth_content}
    \begin{tabular}{lcc}
        \toprule
        Quantity & Single MC & Double MC \\
        \midrule
        Mean truth photons per event & 0.493 & 0.495 \\
        Events with $\geq 2$ denom.\ photons & 0.22\% & 0.21\% \\
        Mean truth \isoETtruth [\GeV] & 0.17 & 0.26 \\
        \bottomrule
    \end{tabular}
\end{table}

The dominant $88\%$ matching-artifact component is entirely driven by
vertex-shifted kinematics (a geometric reconstruction effect, not a
physics effect), and is removed by switching to track-ID matching. The
residual $12\%$ is genuine cluster-loss and is kept as a
genuine pileup degradation in the reconstruction-stage efficiency.

\subsection{Baseline Rejection on Single MC (\DR Cut is Never Benign)}
\label{subsec:single_mc_baseline}

A secondary question is whether $\DR<0.1$ matters outside of pileup
at all. It does: even in single-interaction MC, $\DR<0.1$ rejects
$\sim\!3.5\%$ of correctly \texttt{trkID}-matched truth photons
(Table~\ref{tab:dr_single}), and the rejection is driven entirely by
vertex mismatch ($|\Delta z|>9.35$\,cm) rather than cluster quality:
73\% of rejected clusters have $\ET/\pT^{\text{truth}}\in[0.8,1.2]$,
i.e.\ a nearly-correct reconstructed energy. Six physics arguments
support the removal --- the track-ID is ground truth; the $\DR$ cut
tests vertex resolution rather than calorimeter performance; the
response matrix already absorbs vertex-induced kinematic shifts; and
the $\sim\!3\%$ bias is $z$-dependent and therefore propagates into
any period-dependent quantity. Track-ID-only matching is therefore
adopted as the nominal both for the efficiency and for the DI-MC
study.

\begin{table}[H]
    \centering
    \caption{$\DR<0.1$ failure rates in single-interaction MC.
      ``Fail $\DR$'' is the fraction of \texttt{trkID}-matched photons
      rejected by the geometric cut.}
    \label{tab:dr_single}
    \begin{tabular}{lccc}
        \toprule
        Sample & Truth photons & No trkID match & Fail $\DR<0.1$ (of matched) \\
        \midrule
        photon5  & 273{,}017     & 3.42\% & 2.70\% \\
        photon10 & 2{,}380{,}500 & 2.37\% & 2.98\% \\
        photon20 & 2{,}840{,}703 & 1.72\% & 4.01\% \\
        \midrule
        Combined & 5{,}494{,}220 & 2.09\% & 3.50\% \\
        \bottomrule
    \end{tabular}
\end{table}

\subsection{Per-Stage Efficiency Under Track-ID Matching}
\label{subsec:per_stage_eff}

With $\DR<0.1$ removed, the factorized efficiency chain
$\varepsilon_{\text{reco}} \to \varepsilon_{\text{iso}|\text{reco}} \to
\varepsilon_{\text{id}|\text{reco+iso}} \to \varepsilon_{\text{all}}$
reveals the genuine pileup effects stage by stage.
Table~\ref{tab:eff_breakdown} presents the integrated efficiency for
single, double, and mixed (0\,mrad) photon10 MC under
\texttt{trkID}-only matching; Figure~\ref{fig:eff_overlay_4panel}
shows the same information as a function of truth \pT.

\begin{table}[H]
    \centering
    \caption{Per-component efficiency with \texttt{trkID}-only matching
      (no $\DR$ cut), photon10 samples integrated over
      $\etg = 14$--$30$\,\GeV. Mixed = $0.776 \times \text{single} +
      0.224 \times \text{double}$ (0\,mrad cluster-weighted fractions
      from Table~\ref{tab:pileup_fractions}).}
    \label{tab:eff_breakdown}
    \begin{tabular}{lcccc}
        \toprule
        Efficiency component & Single MC & Double MC & Mixed 0~mrad & Effect (single $\to$ double) \\
        \midrule
        $\varepsilon_{\text{reco}}$ (reconstruction)
          & 93.5\% & 73.1\% & 88.4\% & $-20$~pp (cluster loss + shifted kinematics) \\
        $\varepsilon_{\text{iso}}$ (isolation $|$ reco)
          & 74.2\% & 71.0\% & 73.5\% & $-3$~pp (extra cone energy) \\
        $\varepsilon_{\text{id}}$ (BDT $|$ reco+iso)
          & 82.7\% & 69.4\% & 80.0\% & $\mathbf{-13}$~\textbf{pp} (vertex-shifted kinematics) \\
        \midrule
        $\varepsilon_{\text{all}}$ (total)
          & 57.4\% & 36.0\% & 52.0\% & $-21$~pp (reco + BDT dominate) \\
        \bottomrule
    \end{tabular}
\end{table}

Three structural observations:

\begin{itemize}
\item \textbf{Reconstruction} ($\varepsilon_{\text{reco}}$): the
  20\,pp gap between single (93.5\%) and double (73.1\%) MC splits into
  genuine cluster loss (no-\texttt{trkID}-match rate: 7.9\% double
  vs.\ 2.4\% single, i.e.\ $\sim\!5.5$\,pp of the drop) and
  common-cut failures from shifted kinematics
  ($\sim\!15$\,pp).
\item \textbf{Isolation} ($\varepsilon_{\text{iso}}$): essentially
  unchanged (3\,pp degradation). Extra cone energy from the second
  collision is present but small compared with the 4\,\GeV isolation
  threshold.
\item \textbf{Photon~ID} ($\varepsilon_{\text{id}}$): a \textbf{13\,pp
  degradation} in double MC (69.4\% vs.\ 82.7\%), entirely driven by
  vertex-displaced shower-shape inputs to the parametric
  $\ET$-dependent BDT threshold. In single MC the same effect is
  negligible ($-0.2$\,pp) because only $\sim\!2.3\%$ of clusters have
  large vertex shifts. In double MC, $\sim\!58\%$ of recovered clusters
  are ``new'' (had $\DR>0.1$), and these have a BDT pass rate of only
  59.8\% vs.\ 82.7\% for the full population.
\end{itemize}

\begin{figure}[H]
    \centering
    \includegraphics[width=0.95\textwidth]{../plotting/figures/efficiency_overlay/eff_overlay_noDR_4panel.pdf}
    \caption{Per-stage efficiency vs.\ truth \pT without the $\DR$ cut,
      single (blue), double (red), and mixed 0\,mrad (black) photon10
      MC. Upper left: reconstruction. Upper right: isolation.
      Lower left: photon ID (BDT). Lower right: total. The BDT stage
      shows the cleanest 13\,pp gap between single and double MC, driven
      by vertex-shifted kinematics feeding an $\ET$-dependent BDT
      threshold.}
    \label{fig:eff_overlay_4panel}
\end{figure}

\begin{figure}[H]
    \centering
    \includegraphics[width=0.75\textwidth]{../plotting/figures/efficiency_overlay/eff_id_newclusters_overlay.pdf}
    \caption{BDT pass rate for ``new'' clusters (those with $\DR>0.1$
      but valid \texttt{trkID}) vs.\ truth \pT. These vertex-shifted
      clusters have 10--13\,pp lower BDT pass rates than the full
      population, with the deficit narrowing at higher \pT where the
      parametric BDT score is further above threshold.}
    \label{fig:eff_id_newclusters}
\end{figure}

For reference, Table~\ref{tab:eff_withdr} shows the legacy efficiency
with $\DR<0.1$ retained. The large reco drop there is the matching
artifact; the BDT appears ``unaffected'' only because vertex-shifted
clusters have already been pre-filtered out of the denominator.

\begin{table}[H]
    \centering
    \caption{Per-component efficiency under the legacy $\DR<0.1$ match,
      for reference. The reco drop is the matching artifact; the BDT
      only looks unaffected because vertex-shifted clusters have been
      silently rejected upstream.}
    \label{tab:eff_withdr}
    \begin{tabular}{lccc}
        \toprule
        Efficiency component & Single MC & Double MC & Mixed 0~mrad \\
        \midrule
        $\varepsilon_{\text{reco}}$ ($\DR<0.1$) & 91.4\% & 30.7\% & 76.2\% \\
        $\varepsilon_{\text{iso}}$ & 74.2\% & 71.1\% & 73.9\% \\
        $\varepsilon_{\text{id}}$ (BDT) & 82.9\% & 82.6\% & 82.9\% \\
        \midrule
        $\varepsilon_{\text{all}}$ & 56.2\% & 18.0\% & 46.6\% \\
        \bottomrule
    \end{tabular}
\end{table}

\subsection{\texttt{TEfficiency} Weighted-Event Bug Fix}
\label{subsec:tefficiency_fix}

During the $\DR$ investigation a bug was identified in the
\texttt{TEfficiency} usage in \texttt{RecoEffCalculator\_TTreeReader.C}:
the previous \texttt{SetWeight(w) + Fill(pass, pT)} pattern stored
unweighted event counts, and \texttt{hadd}-ing across samples would
therefore combine them by raw event count rather than by cross-section
weight. The corrected implementation is
\texttt{SetUseWeightedEvents()} plus \texttt{FillWeighted(pass, weight,
pT)}, which stores weighted sums and composes correctly under
\texttt{hadd}. All numbers in Table~\ref{tab:eff_breakdown} use the
corrected implementation.

\subsection{Implications for the Nominal Analysis}
\label{subsec:eff_implications}

\begin{itemize}
\item \textbf{Reco loss contributes $1.6\%$ to the 1.5\,mrad
  cross section}: 20\,pp single$\to$double efficiency drop times
  the 7.9\% cluster-weighted pileup fraction.
\item \textbf{BDT degradation contributes $1.0\%$}:
  $0.079 \times 13\,\text{pp}$.
\item \textbf{Isolation contributes $0.2\%$}:
  $0.079 \times 3\,\text{pp}$.
\item \textbf{Total integrated DI effect on the 1.5\,mrad cross
  section}: $1.6\%+1.0\%+0.2\%=\mathbf{2.8\%}$ pre-correction,
  partially correlated with the vertex-shift systematic that is
  already evaluated separately. At the total-efficiency level
  (mixed 0\,mrad), the same decomposition gives
  $\varepsilon_\mathrm{all}^\mathrm{mix}/\varepsilon_\mathrm{all}^\mathrm{single}
  = 52.0/57.4$, i.e.\ a $1-52.0/57.4=\mathbf{9.4\%}$ relative
  reduction (Table~\ref{tab:eff_breakdown}). The 1.5\,mrad effect
  is smaller than this because the pileup fraction is $\sim\!3\times$
  smaller than at 0\,mrad.
\end{itemize}

The nominal analysis uses \textbf{single-interaction MC for the
efficiency calculation} (with \texttt{trkID}-only matching). Mixed MC
is reserved for the shower-shape data--MC validation
(\S\ref{sec:shower_shape}) and the ABCD background-template
cross-check. Double MC is not used as a default because the apparent
bulk efficiency drop is dominated by the matching artifact rather than
a genuine physics effect: closing it would require simulating the
low-pileup target precisely rather than averaging a broad
generator profile.


%% ==================================================================
\section{Shower-Shape Data-vs-Mixed-MC Validation}
\label{sec:shower_shape}

\subsection{Two-Pass Blending Pipeline}
\label{subsec:two_pass}

The pileup blend is built with \texttt{run\_showershape\_double.sh}
in two passes over \texttt{ShowerShapeCheck.C}. The first pass
computes a blended MC vertex-$z$ template, the second runs the full
shower-shape analysis using it:

\paragraph{Pass~1 (vertex scan).}
\texttt{ShowerShapeCheck.C} is called with
\texttt{do\_vertex\_scan=true} and fills only the \zvtx histogram,
weighting each sample by its cluster-level fraction
($w_{\text{single}}$ or $w_{\text{double}}$ from
Table~\ref{tab:pileup_fractions}). Five parallel jobs run:
\texttt{photon10\_double(0.224)}, \texttt{jet12\_double(0.224)},
\texttt{photon10\_nom(0.776)}, \texttt{jet12\_nom(0.776)}, and data.
The output vtxscan ROOT files are merged per-period via \texttt{hadd}
to form the blended MC vertex distribution.

\paragraph{Pass~2 (full analysis).}
\texttt{ShowerShapeCheck.C} is re-run with \texttt{mix\_weight}
(multiplicative fraction applied to all histogram fills) and
\texttt{vtxscan\_sim\_override} (path to the blended vtxscan file
from Pass~1) set. All MC samples use the \emph{same} blended vtxscan
for vertex reweighting, which ensures consistent \zvtx templates
across samples regardless of which sample originally provided the
vertex distribution. Each sample's fills are scaled by its mixing
weight on top of the cross-section weight. Single and double outputs
are \texttt{hadd}-merged to produce the final blended MC.

The full signature is
\begin{quote}\footnotesize\ttfamily
ShowerShapeCheck(configname, filetype, doinclusive, do\_vertex\_scan, mix\_weight, vtxscan\_sim\_override)
\end{quote}
with \texttt{mix\_weight} = 0.224 (or 0.776) at 0\,mrad and 0.079
(or 0.921) at 1.5\,mrad, and \texttt{vtxscan\_sim\_override} pointing
at the Pass-1 blended ROOT file. The same script accepts a single
command-line argument overriding the default 0\,mrad fraction, so
running the 1.5\,mrad blend is simply \texttt{bash
run\_showershape\_double.sh config\_showershape.yaml 0.079}.

\paragraph{Truth-vertex reweighting.} Starting from the April 15
regeneration, the shower-shape blending pipeline applies the iterative
truth-vertex reweight (\S\ref{sec:truth_vertex}) via the
\texttt{truth\_vertex\_reweight\_on} / \texttt{truth\_vertex\_reweight\_file}
fields in the showershape configs, rather than the reco-level
$f(z_r)$ alone. The samples shown in this section use this nominal
setting.

\subsection{Comparison Methodology}
\label{subsec:shower_method}

Shower-shape distributions are compared at two selection levels:
\begin{itemize}
\item \textbf{cut0}: no selection. Data includes NPB contamination
  that is not modeled in MC, which complicates interpretation.
\item \textbf{cut1}: after the NPB preselection
  (NPB score $>0.5$). This removes non-physics clusters and provides
  the cleanest pileup assessment. \textbf{All plots in this section
  use cut1.}
\end{itemize}
Each comparison plot shows data (black points), signal MC (red),
and inclusive MC (blue), area-normalized, with a \chidf and $p$-value
annotation and a Data$\,-\,$Inclusive~MC difference panel. Left
panels show nominal (single-only) MC; right panels show mixed
(double-blended) MC.

\subsection{0\,mrad Period (22.4\% Double Interaction)}
\label{subsec:results_0mrad}

At 0\,mrad the pileup fraction is high and the nominal (single-only)
MC misrepresents the data at the cut0/cut1 shower-shape level for
pileup-sensitive variables. Mixing in 22.4\% of DI MC closes the
gap substantially for \wetacog at $\ET\geq 14$\,\GeV.

\subsubsection{\wetacog (eta cluster width)}

\wetacog is the variable most sensitive to DI: vertex shifts along
the beam axis broaden the reconstructed $\eta$ profile directly via
Eq.~\eqref{eq:deta_vs_dz}.

\begin{figure}[H]
\centering
\begin{subfigure}[t]{0.48\textwidth}
\includegraphics[width=\textwidth]{../plotting/figures/showershape_0rad/dis_weta_cogx_eta0_pt0_cut1.pdf}
\caption{Nominal MC}
\end{subfigure}\hfill
\begin{subfigure}[t]{0.48\textwidth}
\includegraphics[width=\textwidth]{../plotting/figures/showershape_0rad/dis_mixed_weta_cogx_eta0_pt0_cut1.pdf}
\caption{Mixed MC (22.4\% DI)}
\end{subfigure}
\caption{\wetacog, $10<\ET<14$\,\GeV (pt0), 0\,mrad, cut1.
  Nominal \chidf = 47.40; mixed \chidf = 28.96 ($1.6\times$
  improvement).}
\label{fig:shower_0mrad_weta_pt0}
\end{figure}

\begin{figure}[H]
\centering
\begin{subfigure}[t]{0.48\textwidth}
\includegraphics[width=\textwidth]{../plotting/figures/showershape_0rad/dis_weta_cogx_eta0_pt1_cut1.pdf}
\caption{Nominal MC}
\end{subfigure}\hfill
\begin{subfigure}[t]{0.48\textwidth}
\includegraphics[width=\textwidth]{../plotting/figures/showershape_0rad/dis_mixed_weta_cogx_eta0_pt1_cut1.pdf}
\caption{Mixed MC (22.4\% DI)}
\end{subfigure}
\caption{\wetacog, $14<\ET<18$\,\GeV (pt1), 0\,mrad, cut1.
  Nominal \chidf = 17.64; mixed \chidf = 2.99 ($5.9\times$
  improvement).}
\label{fig:shower_0mrad_weta_pt1}
\end{figure}

\begin{figure}[H]
\centering
\begin{subfigure}[t]{0.48\textwidth}
\includegraphics[width=\textwidth]{../plotting/figures/showershape_0rad/dis_weta_cogx_eta0_pt3_cut1.pdf}
\caption{Nominal MC}
\end{subfigure}\hfill
\begin{subfigure}[t]{0.48\textwidth}
\includegraphics[width=\textwidth]{../plotting/figures/showershape_0rad/dis_mixed_weta_cogx_eta0_pt3_cut1.pdf}
\caption{Mixed MC (22.4\% DI)}
\end{subfigure}
\caption{\wetacog, $22<\ET<28$\,\GeV (pt3), 0\,mrad, cut1.
  Nominal \chidf = 15.58; mixed \chidf = 0.75 ($p=0.84$;
  $20.8\times$ improvement). The mixed MC achieves $p=0.84$, i.e.\
  the data-MC residual in this bin is consistent with statistical
  fluctuation.}
\label{fig:shower_0mrad_weta_pt3}
\end{figure}

The improvement grows with \ET: at $22$--$28$\,\GeV the nominal MC
discrepancy becomes large relative to statistical uncertainties, and
the mixed MC provides the largest relative correction.

\subsubsection{BDT score}

The photon/jet BDT amplifies underlying shower-shape distortions
through its nonlinear decision function and is the most sensitive
composite variable at the cluster level.

\begin{figure}[H]
\centering
\begin{subfigure}[t]{0.48\textwidth}
\includegraphics[width=\textwidth]{../plotting/figures/showershape_0rad/dis_bdt_eta0_pt0_cut1.pdf}
\caption{Nominal MC}
\end{subfigure}\hfill
\begin{subfigure}[t]{0.48\textwidth}
\includegraphics[width=\textwidth]{../plotting/figures/showershape_0rad/dis_mixed_bdt_eta0_pt0_cut1.pdf}
\caption{Mixed MC (22.4\% DI)}
\end{subfigure}
\caption{BDT score, $10<\ET<14$\,\GeV (pt0), 0\,mrad, cut1.
  Nominal \chidf = 6.85; mixed \chidf = \textcolor{red}{20.30}
  ($3.0\times$ \textbf{degradation}). See
  \S\ref{subsec:lowet_discussion} for the discussion of the low-\ET
  crossover.}
\label{fig:shower_0mrad_bdt_pt0}
\end{figure}

\begin{figure}[H]
\centering
\begin{subfigure}[t]{0.48\textwidth}
\includegraphics[width=\textwidth]{../plotting/figures/showershape_0rad/dis_bdt_eta0_pt1_cut1.pdf}
\caption{Nominal MC}
\end{subfigure}\hfill
\begin{subfigure}[t]{0.48\textwidth}
\includegraphics[width=\textwidth]{../plotting/figures/showershape_0rad/dis_mixed_bdt_eta0_pt1_cut1.pdf}
\caption{Mixed MC (22.4\% DI)}
\end{subfigure}
\caption{BDT score, $14<\ET<18$\,\GeV (pt1), 0\,mrad, cut1.
  Nominal \chidf = 12.20; mixed \chidf = 1.63 ($p=0.004$;
  $7.5\times$ improvement).}
\label{fig:shower_0mrad_bdt_pt1}
\end{figure}

\begin{figure}[H]
\centering
\begin{subfigure}[t]{0.48\textwidth}
\includegraphics[width=\textwidth]{../plotting/figures/showershape_0rad/dis_bdt_eta0_pt3_cut1.pdf}
\caption{Nominal MC}
\end{subfigure}\hfill
\begin{subfigure}[t]{0.48\textwidth}
\includegraphics[width=\textwidth]{../plotting/figures/showershape_0rad/dis_mixed_bdt_eta0_pt3_cut1.pdf}
\caption{Mixed MC (22.4\% DI)}
\end{subfigure}
\caption{BDT score, $22<\ET<28$\,\GeV (pt3), 0\,mrad, cut1.
  Nominal \chidf = 9.62; mixed \chidf = 1.75 ($p=0.001$;
  $5.5\times$ improvement).}
\label{fig:shower_0mrad_bdt_pt3}
\end{figure}

\subsubsection{$f_1$ and $E_{1\times1}/E_{3\times3}$}

Double interactions spread energy outward from the shower core,
reducing the inner-ring energy fraction $f_1$. The
$E_{1\times1}/E_{3\times3}$ ratio is less sensitive to pileup
broadening.

\begin{figure}[H]
\centering
\begin{subfigure}[t]{0.48\textwidth}
\includegraphics[width=\textwidth]{../plotting/figures/showershape_0rad/dis_et1_eta0_pt1_cut1.pdf}
\caption{$f_1$, nominal MC}
\end{subfigure}\hfill
\begin{subfigure}[t]{0.48\textwidth}
\includegraphics[width=\textwidth]{../plotting/figures/showershape_0rad/dis_mixed_et1_eta0_pt1_cut1.pdf}
\caption{$f_1$, mixed MC (22.4\% DI)}
\end{subfigure}
\\[6pt]
\begin{subfigure}[t]{0.48\textwidth}
\includegraphics[width=\textwidth]{../plotting/figures/showershape_0rad/dis_e11_to_e33_eta0_pt1_cut1.pdf}
\caption{$E_{1\times1}/E_{3\times3}$, nominal MC}
\end{subfigure}\hfill
\begin{subfigure}[t]{0.48\textwidth}
\includegraphics[width=\textwidth]{../plotting/figures/showershape_0rad/dis_mixed_e11_to_e33_eta0_pt1_cut1.pdf}
\caption{$E_{1\times1}/E_{3\times3}$, mixed MC (22.4\% DI)}
\end{subfigure}
\caption{$f_1$ (top) and $E_{1\times1}/E_{3\times3}$ (bottom) at
  $14<\ET<18$\,\GeV (pt1), 0\,mrad, cut1. $f_1$: \chidf =
  9.33 $\to$ 1.92 ($4.9\times$). $E_{1\times1}/E_{3\times3}$:
  \chidf = 4.50 $\to$ 4.14 ($1.1\times$).}
\label{fig:shower_0mrad_et1_e11}
\end{figure}

\subsection{1.5\,mrad Period (7.9\% Double Interaction)}
\label{subsec:results_1p5mrad}

At 1.5\,mrad the DI fraction is $\sim\!3\times$ smaller than at
0\,mrad, so the mixed MC correction is a few times smaller. The
nominal agreement at cut1 is already good for most variables
(\chidf $\lesssim 10$); where mixed MC helps, the improvement is
mainly at the high-\ET bins.

\subsubsection{\wetacog}

\begin{figure}[H]
\centering
\begin{subfigure}[t]{0.48\textwidth}
\includegraphics[width=\textwidth]{../plotting/figures/showershape_1p5rad/dis_weta_cogx_eta0_pt0_cut1.pdf}
\caption{Nominal MC}
\end{subfigure}\hfill
\begin{subfigure}[t]{0.48\textwidth}
\includegraphics[width=\textwidth]{../plotting/figures/showershape_1p5rad/dis_mixed_weta_cogx_eta0_pt0_cut1.pdf}
\caption{Mixed MC (7.9\% DI)}
\end{subfigure}
\caption{\wetacog, $10<\ET<14$\,\GeV (pt0), 1.5\,mrad, cut1.
  Nominal \chidf = 9.72; mixed \chidf = \textcolor{red}{42.31}
  ($4.4\times$ \textbf{degradation}).}
\label{fig:shower_1p5_weta_pt0}
\end{figure}

\begin{figure}[H]
\centering
\begin{subfigure}[t]{0.48\textwidth}
\includegraphics[width=\textwidth]{../plotting/figures/showershape_1p5rad/dis_weta_cogx_eta0_pt1_cut1.pdf}
\caption{Nominal MC}
\end{subfigure}\hfill
\begin{subfigure}[t]{0.48\textwidth}
\includegraphics[width=\textwidth]{../plotting/figures/showershape_1p5rad/dis_mixed_weta_cogx_eta0_pt1_cut1.pdf}
\caption{Mixed MC (7.9\% DI)}
\end{subfigure}
\caption{\wetacog, $14<\ET<18$\,\GeV (pt1), 1.5\,mrad, cut1.
  Nominal \chidf = 6.73; mixed \chidf = 4.86 ($1.4\times$
  improvement).}
\label{fig:shower_1p5_weta_pt1}
\end{figure}

\begin{figure}[H]
\centering
\begin{subfigure}[t]{0.48\textwidth}
\includegraphics[width=\textwidth]{../plotting/figures/showershape_1p5rad/dis_weta_cogx_eta0_pt3_cut1.pdf}
\caption{Nominal MC}
\end{subfigure}\hfill
\begin{subfigure}[t]{0.48\textwidth}
\includegraphics[width=\textwidth]{../plotting/figures/showershape_1p5rad/dis_mixed_weta_cogx_eta0_pt3_cut1.pdf}
\caption{Mixed MC (7.9\% DI)}
\end{subfigure}
\caption{\wetacog, $22<\ET<28$\,\GeV (pt3), 1.5\,mrad, cut1.
  Nominal \chidf = 21.10; mixed \chidf = 1.31 ($p=0.118$;
  $16.1\times$ improvement).}
\label{fig:shower_1p5_weta_pt3}
\end{figure}

\subsubsection{BDT score}

\begin{figure}[H]
\centering
\begin{subfigure}[t]{0.48\textwidth}
\includegraphics[width=\textwidth]{../plotting/figures/showershape_1p5rad/dis_bdt_eta0_pt0_cut1.pdf}
\caption{Nominal MC}
\end{subfigure}\hfill
\begin{subfigure}[t]{0.48\textwidth}
\includegraphics[width=\textwidth]{../plotting/figures/showershape_1p5rad/dis_mixed_bdt_eta0_pt0_cut1.pdf}
\caption{Mixed MC (7.9\% DI)}
\end{subfigure}
\caption{BDT score, $10<\ET<14$\,\GeV (pt0), 1.5\,mrad, cut1.
  Nominal \chidf = 5.98; mixed \chidf = \textcolor{red}{33.08}
  ($5.5\times$ \textbf{degradation}).}
\label{fig:shower_1p5_bdt_pt0}
\end{figure}

\begin{figure}[H]
\centering
\begin{subfigure}[t]{0.48\textwidth}
\includegraphics[width=\textwidth]{../plotting/figures/showershape_1p5rad/dis_bdt_eta0_pt1_cut1.pdf}
\caption{Nominal MC}
\end{subfigure}\hfill
\begin{subfigure}[t]{0.48\textwidth}
\includegraphics[width=\textwidth]{../plotting/figures/showershape_1p5rad/dis_mixed_bdt_eta0_pt1_cut1.pdf}
\caption{Mixed MC (7.9\% DI)}
\end{subfigure}
\caption{BDT score, $14<\ET<18$\,\GeV (pt1), 1.5\,mrad, cut1.
  Nominal \chidf = 1.16 ($p=0.207$); mixed \chidf = \textcolor{red}{2.31}
  ($2.0\times$ \textbf{degradation}; nominal already excellent).}
\label{fig:shower_1p5_bdt_pt1}
\end{figure}

\begin{figure}[H]
\centering
\begin{subfigure}[t]{0.48\textwidth}
\includegraphics[width=\textwidth]{../plotting/figures/showershape_1p5rad/dis_bdt_eta0_pt3_cut1.pdf}
\caption{Nominal MC}
\end{subfigure}\hfill
\begin{subfigure}[t]{0.48\textwidth}
\includegraphics[width=\textwidth]{../plotting/figures/showershape_1p5rad/dis_mixed_bdt_eta0_pt3_cut1.pdf}
\caption{Mixed MC (7.9\% DI)}
\end{subfigure}
\caption{BDT score, $22<\ET<28$\,\GeV (pt3), 1.5\,mrad, cut1.
  Nominal \chidf = 1.89; mixed \chidf = 1.13 ($p=0.249$;
  $1.7\times$ improvement).}
\label{fig:shower_1p5_bdt_pt3}
\end{figure}

\subsubsection{$f_1$ and $E_{1\times1}/E_{3\times3}$}

\begin{figure}[H]
\centering
\begin{subfigure}[t]{0.48\textwidth}
\includegraphics[width=\textwidth]{../plotting/figures/showershape_1p5rad/dis_et1_eta0_pt1_cut1.pdf}
\caption{$f_1$, nominal MC}
\end{subfigure}\hfill
\begin{subfigure}[t]{0.48\textwidth}
\includegraphics[width=\textwidth]{../plotting/figures/showershape_1p5rad/dis_mixed_et1_eta0_pt1_cut1.pdf}
\caption{$f_1$, mixed MC (7.9\% DI)}
\end{subfigure}
\\[6pt]
\begin{subfigure}[t]{0.48\textwidth}
\includegraphics[width=\textwidth]{../plotting/figures/showershape_1p5rad/dis_e11_to_e33_eta0_pt1_cut1.pdf}
\caption{$E_{1\times1}/E_{3\times3}$, nominal MC}
\end{subfigure}\hfill
\begin{subfigure}[t]{0.48\textwidth}
\includegraphics[width=\textwidth]{../plotting/figures/showershape_1p5rad/dis_mixed_e11_to_e33_eta0_pt1_cut1.pdf}
\caption{$E_{1\times1}/E_{3\times3}$, mixed MC (7.9\% DI)}
\end{subfigure}
\caption{$f_1$ (top) and $E_{1\times1}/E_{3\times3}$ (bottom) at
  $14<\ET<18$\,\GeV (pt1), 1.5\,mrad, cut1. $f_1$: \chidf =
  4.17 $\to$ 2.55 ($1.6\times$). $E_{1\times1}/E_{3\times3}$:
  \chidf = 6.88 $\to$ 6.04 ($1.1\times$).}
\label{fig:shower_1p5_et1_e11}
\end{figure}

\subsection{Single-vs-Double Overlay with Data from Both Periods}
\label{subsec:overlay}

Figures~\ref{fig:shower_overlay_weta_bdt} and
\ref{fig:shower_overlay_et1_e11} overlay pure single-interaction MC
(red) with pure double-interaction MC (blue) and data from both periods:
0\,mrad (filled black circles) and 1.5\,mrad (open green circles), all
area-normalized at cut1. Data at 0\,mrad systematically tracks the
double-MC tail while data at 1.5\,mrad tracks closer to single MC,
consistent with the 22.4\% vs.\ 7.9\% DI fractions and providing a
direct sample-level cross-check on the blending pipeline. The 8-panel
overlay is split into two 4-panel figures for readability: the
vertex-shift-sensitive observables (\wetacog and the composite BDT
score) in Fig.~\ref{fig:shower_overlay_weta_bdt}, and the
ring-energy observables ($f_1$ and $E_{1\times1}/E_{3\times3}$) in
Fig.~\ref{fig:shower_overlay_et1_e11}.

\begin{figure}[htbp]
\centering
\begin{subfigure}[t]{0.48\textwidth}
\includegraphics[width=\textwidth]{../plotting/figures/double_interaction/weta_cogx_eta0_pt0_data_incl.pdf}
\caption{\wetacog, $10$--$14$\,\GeV}
\end{subfigure}\hfill
\begin{subfigure}[t]{0.48\textwidth}
\includegraphics[width=\textwidth]{../plotting/figures/double_interaction/weta_cogx_eta0_pt2_data_incl.pdf}
\caption{\wetacog, $18$--$22$\,\GeV}
\end{subfigure}\\[6pt]
\begin{subfigure}[t]{0.48\textwidth}
\includegraphics[width=\textwidth]{../plotting/figures/double_interaction/bdt_score_eta0_pt0_data_incl.pdf}
\caption{BDT score, $10$--$14$\,\GeV}
\end{subfigure}\hfill
\begin{subfigure}[t]{0.48\textwidth}
\includegraphics[width=\textwidth]{../plotting/figures/double_interaction/bdt_score_eta0_pt2_data_incl.pdf}
\caption{BDT score, $18$--$22$\,\GeV}
\end{subfigure}
\caption{Sample-level shower-shape overlay, vertex-shift-sensitive
  observables: single-interaction MC (red), double-interaction MC
  (blue), data 0\,mrad (filled black circles), data 1.5\,mrad (open
  green circles), all area-normalized. Left column:
  $10<\ET<14$\,\GeV. Right column: $18<\ET<22$\,\GeV. Top row:
  \wetacog. Bottom row: BDT score. 0\,mrad data systematically tracks
  the DI tail; 1.5\,mrad data tracks closer to single MC.}
\label{fig:shower_overlay_weta_bdt}
\end{figure}

\begin{figure}[htbp]
\centering
\begin{subfigure}[t]{0.48\textwidth}
\includegraphics[width=\textwidth]{../plotting/figures/double_interaction/et1_eta0_pt0_data_incl.pdf}
\caption{$f_1$, $10$--$14$\,\GeV}
\end{subfigure}\hfill
\begin{subfigure}[t]{0.48\textwidth}
\includegraphics[width=\textwidth]{../plotting/figures/double_interaction/et1_eta0_pt2_data_incl.pdf}
\caption{$f_1$, $18$--$22$\,\GeV}
\end{subfigure}\\[6pt]
\begin{subfigure}[t]{0.48\textwidth}
\includegraphics[width=\textwidth]{../plotting/figures/double_interaction/e11_to_e33_eta0_pt0_data_incl.pdf}
\caption{$E_{1\times1}/E_{3\times3}$, $10$--$14$\,\GeV}
\end{subfigure}\hfill
\begin{subfigure}[t]{0.48\textwidth}
\includegraphics[width=\textwidth]{../plotting/figures/double_interaction/e11_to_e33_eta0_pt2_data_incl.pdf}
\caption{$E_{1\times1}/E_{3\times3}$, $18$--$22$\,\GeV}
\end{subfigure}
\caption{Sample-level shower-shape overlay, ring-energy
  observables: colour coding as in
  Fig.~\ref{fig:shower_overlay_weta_bdt}. Left column:
  $10<\ET<14$\,\GeV. Right column: $18<\ET<22$\,\GeV. Top row:
  $f_1$. Bottom row: $E_{1\times1}/E_{3\times3}$. Both observables
  show smaller DI-induced single/double separation than
  \wetacog/BDT, and the DI-sensitivity ordering between periods
  still holds.}
\label{fig:shower_overlay_et1_e11}
\end{figure}

\subsection{\chidf Summary}
\label{subsec:chi2}

Table~\ref{tab:chi2_full} compiles the Data$\,-\,$Inclusive-MC \chidf
and $p$-values for all variables, \ET bins, and crossing angles.
Entries where the mixed MC \emph{degrades} agreement relative to
nominal are highlighted in \textcolor{red}{red}.

\begin{table}[H]
\centering
\caption{Data$\,-\,$Inclusive-MC \chidf and $p$-values at cut1 for
  nominal (single-only) and mixed (double-blended) MC at both
  crossing angles. Lower \chidf is better; $p \ll 1$ indicates
  significant data--MC disagreement. Red entries flag cases where
  the blend degrades agreement.}
\label{tab:chi2_full}
\footnotesize
\begin{tabular}{ll rrrr rrrr}
\toprule
 & & \multicolumn{4}{c}{\textbf{0\,mrad (22.4\% DI)}}
   & \multicolumn{4}{c}{\textbf{1.5\,mrad (7.9\% DI)}} \\
\cmidrule(lr){3-6} \cmidrule(lr){7-10}
Variable & $E_T$ [\GeV]
  & Nom.\ & $p$ & Mixed & $p$
  & Nom.\ & $p$ & Mixed & $p$ \\
\midrule
\wetacog & 10--14 (pt0)
  & 47.40 & $<$0.001 & 28.96 & $<$0.001
  &  9.72 & $<$0.001 & \textcolor{red}{42.31} & $<$0.001 \\
\wetacog & 14--18 (pt1)
  & 17.64 & $<$0.001 &  2.99 & $<$0.001
  &  6.73 & $<$0.001 &  4.86 & $<$0.001 \\
\wetacog & 22--28 (pt3)
  & 15.58 & $<$0.001 &  0.75 & 0.840
  & 21.10 & $<$0.001 &  1.31 & 0.118 \\
\midrule
BDT score & 10--14 (pt0)
  &  6.85 & $<$0.001 & \textcolor{red}{20.30} & $<$0.001
  &  5.98 & $<$0.001 & \textcolor{red}{33.08} & $<$0.001 \\
BDT score & 14--18 (pt1)
  & 12.20 & $<$0.001 &  1.63 & 0.004
  &  1.16 & 0.207 & \textcolor{red}{2.31} & $<$0.001 \\
BDT score & 22--28 (pt3)
  &  9.62 & $<$0.001 &  1.75 & 0.001
  &  1.89 & $<$0.001 &  1.13 & 0.249 \\
\midrule
$f_1$ & 14--18 (pt1)
  &  9.33 & $<$0.001 &  1.92 & $<$0.001
  &  4.17 & $<$0.001 &  2.55 & $<$0.001 \\
\midrule
$E_{1\times1}/E_{3\times3}$ & 14--18 (pt1)
  &  4.50 & $<$0.001 &  4.14 & $<$0.001
  &  6.88 & $<$0.001 &  6.04 & $<$0.001 \\
\bottomrule
\end{tabular}
\end{table}

\subsection{Low-\ET Crossover and Recommendation}
\label{subsec:lowet_discussion}

A consistent pattern emerges across \emph{both} crossing angles:
the DI blending \textbf{degrades} data--MC agreement for the BDT
score at low \ET (pt0, $10$--$14$\,\GeV) while \textbf{improving} it
at higher \ET. This is now confirmed for 0\,mrad ($3.0\times$ BDT
degradation at pt0) and 1.5\,mrad ($5.5\times$ degradation at pt0).
For \wetacog, the crossover lies between pt0 and pt1: at 0\,mrad,
pt0 \wetacog still improves ($1.6\times$); at 1.5\,mrad it
degrades ($4.4\times$). The crossover pileup fraction for \wetacog
at low \ET therefore lies between 7.9\% and 22.4\%.

Three contributing mechanisms are plausible:
\begin{enumerate}
\item \textbf{Shared MC samples}: the same
  \texttt{photon10\_double} and \texttt{jet12\_double} samples are
  used for both crossing angles; the beam conditions in the MC may
  better match one period than the other, and the mismatch only
  becomes visible at low \ET where the shower-shape distributions
  are narrow.
\item \textbf{Low-\ET sensitivity}: clusters near the 8--14\,\GeV
  threshold are most sensitive to absolute vertex shifts, since
  a fixed $\Delta z$ produces a larger fractional kinematic offset.
\item \textbf{Small-fraction amplification at 1.5\,mrad}: at 7.9\%
  mixing, any residual mismatch between the DI MC and the real
  pileup is amplified relative to the small physical correction.
\end{enumerate}

These observations motivate the use of mixed MC as a shower-shape
\emph{validation} rather than as a shower-shape template correction,
and the recommendation: apply mixed MC as a cross-check at
$\ET\geq 14$\,\GeV (0\,mrad) or $\ET\geq 18$\,\GeV (1.5\,mrad), and
take the nominal-mixed spread as a systematic in the lower bins.
The residual shower-shape mismodelling at low \ET is expected to be
subdominant after the truth-vertex reweight is applied
(\S\ref{sec:truth_vertex}), which shares data with the shower-shape
pipeline and whose downstream application carries the same low-\ET
sensitivity.


%% ==================================================================
\section{Truth-Vertex Reweighting}
\label{sec:truth_vertex}

The shower-shape and efficiency analyses use a single-variable
vertex reweight to close the MC luminous-region shape to data. The
production choice is the reco-level ratio
$f(z_r) = D(z_r)/M_{\mathrm{reco}}(z_r)$, which closes the
reco-vertex distribution exactly. This section shows that $f(z_r)$
is \emph{structurally insufficient} for double-interaction events,
quantifies the resulting cross-section bias at 1.5\,mrad, and
presents the nominal fix: an iterative data-driven truth-vertex
reweight $w(z)$, factorized as $w(z_h)$ for single MC and
$w(z_h)\,w(z_\mathrm{mb})$ for double MC.

\subsection{Why the Production $f(z_r)$ Cannot Close the Truth Vertex}
\label{subsec:problem}

The efficiency denominator in
\texttt{RecoEffCalculator\_TTreeReader.C} applies a reco-vertex cut
$|z_r|<\mathtt{vertex\_cut}$ \emph{and} a truth-vertex cut
$|z_h|<\mathtt{vertex\_cut}$ (both set to 60\,cm in
\texttt{config\_bdt\_nom.yaml}). The reco cut is handled correctly
by $f(z_r)$: the reweighted MC reco distribution matches data
bin-by-bin by construction. The truth cut is \emph{not} handled
correctly for double-interaction events.

Decompose the pair of truth vertices
$(z_{t,1},z_{t,2})\equiv (z_h,z_\mathrm{mb})$ of a
double-interaction event into a symmetric coordinate
$\bar z \equiv (z_h+z_\mathrm{mb})/2$ and an antisymmetric
coordinate $\varepsilon \equiv (z_h-z_\mathrm{mb})/2$. Under the
independent Gaussian sampling of the raw MC, $\bar z$ and
$\varepsilon$ are statistically independent zero-mean Gaussians
with $\sigma_{\bar z} = \sigma_\varepsilon = \sigma_\mathrm{MC}/\sqrt 2$.
The reco vertex is $z_r \approx \bar z$ (up to MBD resolution), so
$f(z_r)$ depends \emph{only} on $\bar z$ and is flat in
$\varepsilon$. After reweighting, $\bar z$ matches the target
Gaussian $\mathcal{N}(0,\sigma_\mathrm{lum}^2/2)$ by construction,
but $\varepsilon$ is unchanged. The primary truth vertex
$z_h = \bar z + \varepsilon$ is therefore a sum of a narrowed and
an un-narrowed Gaussian, with variance
\begin{equation}
\sigma^2(z_h)
  = \tfrac{1}{2}\,\sigma_\mathrm{lum}^2 + \tfrac{1}{2}\,\sigma_\mathrm{MC}^2
  = \tfrac{1}{2}(22.3^2 + 62.7^2)
  \simeq (47.0\,\mathrm{cm})^2.
\label{eq:half_narrow_half_wide}
\end{equation}
We use $\sigma_\mathrm{lum} \simeq \sigma_\mathrm{data\,reco}$ here
(22.3\,cm rather than the 21.48\,cm $\sigma_\mathrm{seed}$ quoted in
\S\ref{subsec:method}); MBD resolution is sub-cm, so the two widths
agree to within 1\,cm at 1.5\,mrad and the narrow-Gaussian
approximation to the luminous region is good at the required
precision. Any principled fix must act on the truth vertices
directly, i.e.\ factorize as $w(z_h)\,w(z_\mathrm{mb})$ rather than
depending only on the reco combination $z_r$.

\subsection{Cross-Section Impact at 1.5\,mrad}
\label{subsec:xsec_impact}

For a Gaussian truth distribution of width $\sigma$, the fraction
inside $|z_h|<60$\,cm is $\mathrm{erf}(60/(\sigma\sqrt 2))$:
\begin{itemize}\setlength\itemsep{1pt}
\item $\sigma = 47$\,cm (what nominal $f(z_r)$ gives for double MC):
      $\mathrm{erf}(60/(47\sqrt 2)) \simeq 0.80$
\item $\sigma = 22$\,cm (target for 1.5\,mrad):
      $\mathrm{erf}(60/(22\sqrt 2)) \simeq 0.99$
\end{itemize}
Double-interaction MC therefore loses $\sim\!19$ percentage points
of truth fiducial acceptance that it should not. At 1.5\,mrad the
mixed blend is
$0.921 \times \text{single} + 0.079 \times \text{double}$, so the
bias on the integrated efficiency denominator is
\[
0.079 \times 0.19 \simeq 1.4\,\% \quad \text{at 1.5\,mrad}.
\]
This translates directly to a $+1.4\%$ bias on the extracted cross
section at 1.5\,mrad (the efficiency denominator is too small, so
the quoted cross section is too large). At 0\,mrad the analogous
bias is $\sim\!0.3\%$ because the raw MC beam profile already sits
close to the wider 0\,mrad luminous region. This is comparable to
other PPG12 systematics and worth fixing.

\subsection{Method: Data-Driven Iterative Truth-Vertex Reweight}
\label{subsec:method}

The method derives a truth-vertex weight function $w(z)$ directly
from the data reco distribution, with no closed-form Gaussian
assumption on the luminous-region shape.

\paragraph{Forward model.}
Let $f_d$ be the cluster-weighted double-interaction fraction
($0.079$ at 1.5\,mrad, $0.224$ at 0\,mrad) and
$f_s = 1-f_d$. Under a candidate weight $w(z)$, the mixed-MC reco
distribution is
\begin{equation}
M(z_r) =
f_s \!\!\sum_{e \in T_s}\! w(z_h^{e})\,\delta(z_r - z_r^e)
+
f_d \!\!\sum_{e \in T_d}\! w(z_h^{e})\,w(z_\mathrm{mb}^{e})\,\delta(z_r - z_r^e),
\label{eq:iter_forward}
\end{equation}
evaluated event-by-event over the single-interaction sample $T_s$
and the double-interaction sample $T_d$. For double-interaction
events the weight enters \emph{twice}, once per truth vertex. This
factorization gives the iteration access to the antisymmetric
coordinate $\varepsilon$ of \S\ref{subsec:problem} that the
production reco-only weight $f(z_r)$ cannot touch.

\paragraph{Binning and closure window.}
Both periods use 80 bins over $[-200,200]$\,cm (bin width 5\,cm),
with identical bin edges for $w(z)$ and for the reco-vertex
histograms. The closure window is auto-derived per period as the
$99$-th percentile of $|z_r^\mathrm{data}|$:
1.5\,mrad uses $z_\mathrm{cut}=83$\,cm (34 bins), 0\,mrad uses
$z_\mathrm{cut}=144$\,cm (58 bins). Outside the window, data
statistics are too low for the iteration to constrain $w(z)$.

\paragraph{Seed.}
$w_0(z)=\mathcal{N}(0,\sigma_\mathrm{seed}^2)/\mathcal{N}(0,\sigma_\mathrm{gen}^2)$,
with $\sigma_\mathrm{gen} = 64.997\pm 0.015$\,cm measured from
the photon10 primary truth vertex (identical between periods, since
the same MC is used for both), and
$\sigma_\mathrm{seed}=\mathrm{std}(z_r^\mathrm{data})$ measured
from data --- 21.48\,cm at 1.5\,mrad, 54.32\,cm at 0\,mrad.
$\sigma_\mathrm{seed}$ requires \emph{no} physics input: no
Piwinski narrowing factor, no luminous-region calculation; it is
just the empirical width of the reco distribution we are trying to
match.

\paragraph{Iterative update.}
At iteration $i$ the residual is $r_i(z_r)=D(z_r)/M_i(z_r)$. The
update is Lucy--Richardson-like, multiplicative, and damped:
\begin{equation}
w_{i+1}(z) = w_i(z) \cdot \bigl[\mathrm{smooth}(r_i(z))\bigr]^{\alpha},
\qquad \alpha = 0.5,
\end{equation}
where the residual at reco bin $z_r$ is applied to the weight at
the matching truth bin $z=z_r$. After each update $w$ is
renormalized so $\langle w(z_h)\rangle=1$ averaged over the
single-MC truth distribution (preserving the physics DI fraction).
Damping prevents over-correction; 1-bin Gaussian smoothing on the
residual prevents bin-by-bin oscillation; the update factor is
clipped to $[0.5,2]$ per iteration; boundary bins are edge-clamped.
The final $w(z)$ is data-driven everywhere, without a
post-iteration snap to a seed Gaussian (an earlier variant did
snap; \texttt{compare\_snap\_strategies.py} showed removing the
snap improves 0\,mrad $\chi^2$ by $\sim\!1.3\times$ and
$\max|r-1|$ by $\sim\!3\times$ while leaving 1.5\,mrad essentially
unchanged).

Convergence is $\max|r_i(z)-1|<0.01$ inside the closure window, with
a hard cap of 30 iterations.

\paragraph{Implementation.}
Standalone Python tool in
\texttt{efficiencytool/truth\_vertex\_reweight/}. The mixed-MC
histogram is heavy to rebuild ($\sim\!22$\,M events at 1.5\,mrad,
$\sim\!48$\,M at 0\,mrad), so per-sample arrays are cached to
\texttt{output/cache/*.npz} keyed on file mtime. Each iteration
runs in seconds per iteration on a single core. Outputs live at
\texttt{output/\{1p5mrad,0mrad\}/reweight.root} (containing
\texttt{h\_w\_iterative}, \texttt{h\_w\_seed}, and
\texttt{h\_M\_iterative}) and \texttt{fit\_result.yaml}. The code
is driven from the period configs \texttt{config\_showershape\_\{1p5rad,0rad\}.yaml}
via the new \texttt{truth\_vertex\_reweight\_on} and
\texttt{truth\_vertex\_reweight\_file} fields.

\subsection{Per-Period Results}
\label{subsec:iter_results}

Both periods are fit on the photon10/photon10\_double mix only at
this stage; the jet12 extension is a follow-up
(\S\ref{subsec:iter_caveats}). Table~\ref{tab:iter_summary}
collects the headline numbers.

\begin{table}[H]
\centering
\renewcommand{\arraystretch}{1.15}
\begin{tabular}{l c c}
\toprule
quantity & 1.5\,mrad & 0\,mrad \\
\midrule
run range                                & 51274--54000 & 47289--51274 \\
luminosity [pb$^{-1}$]                   & 16.8588      & 32.6574      \\
$f_\mathrm{double}$ (cluster-weighted)   & 0.079        & 0.224        \\
data $|z_r|$ 99\% percentile [cm]        & 82.0         & 143.9        \\
closure window $z_\mathrm{cut}$ [cm]     & 83           & 144          \\
bins inside closure window               & 34           & 58           \\
$\sigma_\mathrm{seed}=\mathrm{std}(z_r^\mathrm{data})$ [cm]  & 21.48 & 54.32 \\
$\sigma_\mathrm{gen}$ from MC $z_h$ (generator) [cm] & \multicolumn{2}{c}{$64.997 \pm 0.015$} \\
\midrule
initial $\chi^2$ (in window)             & $1{,}419{,}454$ & $26{,}589$ \\
final   $\chi^2$ (in window)             & $453$           & $158$ \\
$\chi^2$ reduction factor                & $3132\times$    & $168\times$ \\
final \chidf (in window)                 & $13.7$ (ndf 33) & $2.78$ (ndf 57) \\
initial $\max|r-1|$                      & $0.635$         & $0.079$ \\
final   $\max|r-1|$ (in window)          & $0.111$         & $0.008$ \\
iterations used                          & 30 (cap hit)    & 8 \\
converged to $0.01$ tolerance            & \textbf{no}     & \textbf{yes} \\
\midrule
(context) parametric Gaussian \chidf     & 3237            & 454 \\
\bottomrule
\end{tabular}
\caption{Iterative data-driven truth-vertex reweight, photon10 mix only,
  with the post-iteration snap to seed removed. The 0\,mrad fit
  converges to the 1\% per-bin tolerance in 8 iterations (final
  $\max|r-1|\simeq 0.8\%$). The 1.5\,mrad fit reduces in-window
  $\chi^2$ by a factor $\sim\!3100$ in 30 iterations and stalls at
  $\max|r-1|\simeq 11\%$, dominated by bins at the closure-window
  edge $|z_r|\simeq 83$\,cm. Inside the bulk $|z_r|<30$\,cm the
  per-bin closure is at the $1$--$2\%$ level. The last row reports
  the parametric best-fit Gaussian \chidf from a 1-parameter
  $\sigma_\mathrm{tgt}$ scan; the iterative method improves on this
  by a factor $\sim\!236$ at 1.5\,mrad and $\sim\!163$ at 0\,mrad.}
\label{tab:iter_summary}
\end{table}

\paragraph{Hero figure.}
Figure~\ref{fig:iter_two_period} is the single-figure summary of
the iterative method across both periods: data reco vs.\ reweighted
MC reco, residuals, the recovered $w(z)$, and the headline
numbers.

\begin{figure}[H]
\centering
\includegraphics[width=\linewidth]{figures/vertex_reweight_iter_two_period.pdf}
\caption{Iterative data-driven truth-vertex reweight summary, both
  crossing angles. Top-left: data reco $z_r$ (solid) vs.\ mixed MC
  reco $z_r$ after iterative reweight (dashed), both periods
  overlaid (log scale). Top-right: residual
  $r(z_r)=D(z_r)/M_\mathrm{iter}(z_r)$ inside each period's closure
  window (vertical dashed lines). Bottom-left: recovered
  truth-vertex weight $w(z)$ for both periods (log scale).
  Bottom-right: summary statistics. The 0\,mrad case is essentially
  closed; the 1.5\,mrad case is closed in the bulk and stalls at
  the closure-window edge.}
\label{fig:iter_two_period}
\end{figure}

\paragraph{Truth vertex comparison across methods.}
Figure~\ref{fig:truth_comparison} overlays the primary truth
vertex $z_h$ distribution of the mixed MC under three weighting
methods: no reweight (raw MC, $\sigma \simeq 64$\,cm), the
production reco-level weight $f(z_r)$, and the new iterative
truth-level weight $w(z_h)$. At 1.5\,mrad the production method
leaves $z_h$ at $\sigma \simeq 27$\,cm --- narrower than the raw
MC but still far wider than the true luminous-region width
($\sigma_\mathrm{lum} \simeq 22$\,cm). The iterative method narrows
$z_h$ to $\sigma \simeq 16$\,cm, which is narrower than the data
reco width $\sigma_\mathrm{seed}=21.5$\,cm precisely as expected
(the MBD vertex finder adds resolution smearing, so each truth
vertex must be narrower than the reco width). At 0\,mrad both
methods converge at $\sigma \simeq 54$--$57$\,cm.

\begin{figure}[H]
\centering
\includegraphics[width=\linewidth]{figures/vertex_reweight_truth_comparison.pdf}
\caption{Truth vertex $z_h$ distribution of the mixed photon10 MC
  under three reweighting methods, per period. Left: 1.5\,mrad,
  where $f(z_r)$ (blue dashed) fails to narrow the truth
  distribution sufficiently while the iterative $w(z_h)$ (red)
  narrows it past the data reco width. Right: 0\,mrad, where both
  methods agree because the MC beam profile is already close to the
  data luminous region.}
\label{fig:truth_comparison}
\end{figure}

\paragraph{0\,mrad: clean convergence.}
Figure~\ref{fig:iter_0mrad_closure} shows the 6-panel closure
diagnostic for 0\,mrad. The seed Gaussian with
$\sigma_\mathrm{seed}\simeq 54$\,cm already captures most of the
shape (initial $\max|r-1|=8\%$, initial \chidf\,$\simeq\!467$);
the parametric best-fit Gaussian at $\sigma_\mathrm{tgt}\simeq 54.6$\,cm
is essentially identical (\chidf\,$\simeq\!454$) because the data
is nearly Gaussian. The iteration cleans up the residual over 8
updates and drops below the 1\%-per-bin tolerance, with final
$\max|r-1|\simeq 0.8\%$ and \chidf\,$\simeq\!2.78$. The recovered
$w(z)$ is smooth and broadly Gaussian with a small central
enhancement.

\begin{figure}[H]
\centering
\includegraphics[width=\linewidth]{figures/vertex_reweight_iter_0mrad_closure.pdf}
\caption{0\,mrad iterative closure (6 panels). Top-left/centre:
  data reco vs.\ mixed MC reco under no reweight (gray dotted),
  the seed Gaussian
  $w(z)\propto\mathcal{N}(0,\sigma_\mathrm{seed}^2)/\mathcal{N}(0,\sigma_\mathrm{gen}^2)$
  (orange dash-dot), the parametric best-fit Gaussian (blue
  dashed), and the iterative reweight (red). Top-right: the
  corresponding $w(z)$. Bottom-left: residual $D/M$ in the closure
  window with Poisson errors and $\pm 1\%$/$\pm 2\%$ reference
  bands. Bottom-right: per-iteration $\chi^2$ and $\max|r-1|$
  trajectory. At 0\,mrad the seed and the parametric Gaussian are
  nearly identical because the data reco is already close to a
  single Gaussian; the iterative method improves on both.}
\label{fig:iter_0mrad_closure}
\end{figure}

\paragraph{1.5\,mrad: bulk closure, edge stall.}
Figure~\ref{fig:iter_1p5mrad_closure} shows the 6-panel diagnostic
for 1.5\,mrad. The seed Gaussian with $\sigma_\mathrm{seed}\simeq 21$\,cm
starts from $\max|r-1|=64\%$ and an in-window $\chi^2$ of
$1.42\times 10^6$. The first $\sim\!5$ iterations make dramatic
progress; the iteration then enters a slow plateau driven by bins
at the closure-window edge $|z_r|\simeq 83$\,cm. After 30
iterations $\chi^2$ in the window has dropped by $\sim\!3100\times$
to $\simeq\!453$, and $\max|r-1|$ sits at $\simeq\!0.11$, set by
those stuck edge bins. Inside the bulk $|z_r|<30$\,cm the closure
is $\sim\!1$--$2\%$ per bin; between 30 and 60\,cm the residual
climbs to $\sim\!5$--$10\%$. The recovered $w(z)$ shows a strong
central core boost and a clear non-Gaussian shape in
$|z|\simeq 40$--$60$\,cm that no symmetric Gaussian ansatz can
represent.

\begin{figure}[H]
\centering
\includegraphics[width=\linewidth]{figures/vertex_reweight_iter_1p5mrad_closure.pdf}
\caption{1.5\,mrad iterative closure (6 panels), same layout as
  Fig.~\ref{fig:iter_0mrad_closure}. At 1.5\,mrad the seed Gaussian
  ($\sigma_\mathrm{seed}=21.5$\,cm, orange dash-dot) and the
  parametric best-fit Gaussian ($\sigma_\mathrm{tgt}=15.7$\,cm,
  blue dashed) are distinct; both mis-predict the reco-vertex
  shape by $30$--$60\%$ at the closure-window edge. The iterative
  method (red points) closes the bulk to $1$--$2\%$ inside
  $|z_r|<30$\,cm and stalls at $\max|r-1|\simeq 0.11$ at
  $|z_r|\simeq 83$\,cm. The recovered $w(z)$ in the top-right
  panel shows a strong central core boost and a clear non-Gaussian
  shape in $|z|\simeq 40$--$60$\,cm.}
\label{fig:iter_1p5mrad_closure}
\end{figure}

\paragraph{Method comparison.}
Figure~\ref{fig:iter_1p5mrad_methods} overlays the iterative
reweight against three reference baselines at 1.5\,mrad: no
reweight, the seed Gaussian (initial point of the iteration), and
the parametric best-fit Gaussian (1-parameter
$\sigma_\mathrm{tgt}$ scan with minimum at 15.7\,cm). The
no-reweight MC overshoots data by a factor of several; both
Gaussian baselines match the width on average but mismatch the
shape at $30$--$70\%$; only the iterative method follows the
narrow data core \emph{and} tracks the non-Gaussian mid-tail.

\begin{figure}[H]
\centering
\includegraphics[width=0.95\linewidth]{figures/vertex_reweight_iter_1p5mrad_methods.pdf}
\caption{Method comparison at 1.5\,mrad: data reco vs.\ mixed MC
  reco under no reweight (gray dotted), seed Gaussian
  $\sigma_\mathrm{seed}=21.5$\,cm (orange dash-dot), parametric
  best-fit Gaussian $\sigma_\mathrm{tgt}=15.7$\,cm (blue dashed),
  and the iterative data-driven reweight (red), with ratio panel
  below. The iterative method is the only one that follows the
  narrow data core while also tracking the non-Gaussian mid-tail
  structure.}
\label{fig:iter_1p5mrad_methods}
\end{figure}

\paragraph{Parametric $\sigma_\mathrm{tgt}$ scan.}
Figure~\ref{fig:sigma_scan} shows the stage-1 parametric $\chi^2$
scan that precedes the iterative method. A single Gaussian
truth-vertex weight
$w(z)\propto\mathcal{N}(0,\sigma_\mathrm{tgt}^2)/\mathcal{N}(0,\sigma_\mathrm{gen}^2)$
is evaluated at each scan point. At 1.5\,mrad the minimum is at
$\sigma_\mathrm{tgt}\simeq 15.7$\,cm with
\chidf\,$\simeq\!3237$; the iterative method reaches 13.7, a
$\sim\!236\times$ improvement. At 0\,mrad the minimum is at
$\sigma_\mathrm{tgt}\simeq 54.6$\,cm with \chidf\,$\simeq\!454$;
the iterative method reaches 2.78, a $\sim\!163\times$
improvement.

\begin{figure}[H]
\centering
\includegraphics[width=0.95\linewidth]{figures/vertex_reweight_sigma_scan.pdf}
\caption{Stage-1 parametric \chidf as a function of the target
  truth-vertex Gaussian width $\sigma_\mathrm{tgt}$, for the
  1.5\,mrad (left) and 0\,mrad (right) periods. Vertical dashed
  line: parabolic-refined minimum. Horizontal dashed line (red):
  final \chidf achieved by the iterative method. Horizontal
  dotted line: \chidf = 1. The gap between the parametric minimum
  and the iterative floor quantifies the non-Gaussian shape
  mismatch that no single-Gaussian ansatz can capture.}
\label{fig:sigma_scan}
\end{figure}

\paragraph{Convergence and pull diagnostics.}
Figures~\ref{fig:iter_convergence_0mrad}--\ref{fig:iter_convergence_1p5mrad}
show the per-iteration trajectory of the $\chi^2$ and
$\max|r-1|$ metrics.
Figures~\ref{fig:iter_pull_0mrad}--\ref{fig:iter_pull_1p5mrad}
show the pull distribution of $D-M$ in the closure window.

\begin{figure}[H]
\centering
\begin{subfigure}[t]{0.48\textwidth}
\includegraphics[width=\textwidth]{figures/vertex_reweight_iter_0mrad_convergence.pdf}
\caption{0\,mrad convergence}
\label{fig:iter_convergence_0mrad}
\end{subfigure}\hfill
\begin{subfigure}[t]{0.48\textwidth}
\includegraphics[width=\textwidth]{figures/vertex_reweight_iter_1p5mrad_convergence.pdf}
\caption{1.5\,mrad convergence}
\label{fig:iter_convergence_1p5mrad}
\end{subfigure}
\caption{Per-iteration $\chi^2$ and $\max|r-1|$ trajectory for
  the iterative reweight. Left: 0\,mrad (converged in 8 iterations).
  Right: 1.5\,mrad (cap hit at 30 iterations, stalled at the
  closure-window edge).}
\end{figure}

\begin{figure}[H]
\centering
\begin{subfigure}[t]{0.48\textwidth}
\includegraphics[width=\textwidth]{figures/vertex_reweight_iter_0mrad_pull.pdf}
\caption{0\,mrad pull}
\label{fig:iter_pull_0mrad}
\end{subfigure}\hfill
\begin{subfigure}[t]{0.48\textwidth}
\includegraphics[width=\textwidth]{figures/vertex_reweight_iter_1p5mrad_pull.pdf}
\caption{1.5\,mrad pull}
\label{fig:iter_pull_1p5mrad}
\end{subfigure}
\caption{Residual pull distributions in the closure window after
  the iterative reweight. 0\,mrad (left) is centred and roughly
  unit-width; 1.5\,mrad (right) shows a small tail at the
  closure-window edge.}
\end{figure}

\subsection{Physics Interpretation}
\label{subsec:iter_interp}

\paragraph{Why 0\,mrad is easy.}
The 0\,mrad luminous region ($\sigma_\mathrm{lum}\simeq 54$\,cm)
is close to the generator beam profile ($\sigma_\mathrm{gen}\simeq 65$\,cm),
and the data reco shape is nearly Gaussian. The seed is already a
good first approximation and the iteration just cleans up small
departures (slightly enhanced central core).

\paragraph{Why 1.5\,mrad is hard.}
Three things conspire:
\begin{itemize}
\item The Piwinski narrowing compresses the luminous region to
  $\sigma_\mathrm{lum}\simeq 22$\,cm, nearly a factor of three
  narrower than the generator profile. The reweight must extract
  a narrow target from a wide source; the dynamic range of $w(z)$
  at the central core is correspondingly large ($w(0)\simeq 8$).
\item The data reco distribution is \emph{not} a clean Gaussian.
  It has a narrow core plus a mid-tail shoulder near
  $|z_r|\simeq 45$--$60$\,cm. The iterative method discovers this
  structure directly --- the recovered $w(z)$ is non-Gaussian in
  the same $z$ range --- whereas a single-parameter Gaussian
  cannot represent it.
\item The MBD resolution tail leak from core-truth events fills
  the bins near the closure-window edge ($|z_r|\simeq 83$\,cm).
  The iteration becomes insensitive to $w(z)$ at that $z$ because
  the contributions to those reco bins come dominantly from truth
  vertices at \emph{other} $z$ values via the wide
  double-interaction reco-truth kernel
  (Appendix~\ref{app:double_diag}). This sets the residual floor
  at $\max|r-1|\simeq 11\%$.
\end{itemize}

\paragraph{What the recovered $w(z)$ looks like.}
At 1.5\,mrad the recovered $w(z)$ has three distinguishing
features (top-right panel of Fig.~\ref{fig:iter_1p5mrad_closure}):
(a) a strong central peak with $w(0)\simeq 8$ that pulls doubles
whose truth vertices both lie near the IP into the data-like core;
(b) a mid-$z$ dip that suppresses the generator's natural
$|z|\sim 30$--$50$\,cm population; and (c) a clear $\pm$ asymmetry
in $|z|\simeq 40$--$60$\,cm tracking the asymmetric mid-tail bump
in the data. None of these features is representable by a
symmetric single-parameter Gaussian. At 0\,mrad, $w(z)$ is
essentially the seed Gaussian with a small central enhancement.

\subsection{Downstream Hookup}
\label{subsec:hookup}

Both \texttt{ShowerShapeCheck.C} and
\texttt{RecoEffCalculator\_TTreeReader.C} now accept two analysis
fields: \texttt{truth\_vertex\_reweight\_on} (0/1) and
\texttt{truth\_vertex\_reweight\_file} (path to \texttt{reweight.root}).
When \texttt{on=1}, the code reads \texttt{h\_w\_iterative},
multiplies the per-event weight by $w(z_h)$ for single-interaction
MC and by $w(z_h)\,w(z_\mathrm{mb})$ for double-interaction MC. The
secondary truth vertex branch is picked up through an
optional-pointer guard, so single-MC trees without the branch fall
through to the single-vertex path. The reco-level and truth-level
reweights are mutually exclusive; the config loader forces
\texttt{vertex\_reweight\_on=0} when
\texttt{truth\_vertex\_reweight\_on=1}. The showershape configs
\texttt{config\_showershape\_\{1p5rad,0rad\}.yaml} are pre-wired
with the matching \texttt{reweight.root}; the nominal BDT configs
carry the same fields but disabled, so flipping the main analysis
over to the truth-vertex reweight is a one-line config edit once
the jet-sample extension is done.

\subsection{Caveats and Next Steps}
\label{subsec:iter_caveats}

\begin{itemize}
\item \textbf{photon10-only mix}: the current fit uses
  \texttt{photon10+photon10\_double} only; the full production
  blend also includes \texttt{jet12+jet12\_double}. Extending the
  fit is mechanically two more \texttt{load\_sample} calls plus
  the appropriate cluster-weighted blend, but the closure numbers
  will move and must be remeasured before any production hookup.
\item \textbf{1.5\,mrad not formally converged}: the 30-iteration
  cap was hit. Open questions: (a) whether extending to
  $\sim\!100$ iterations pushes the stuck edge bin to convergence
  or whether it is a hard physical floor set by the DI
  reco-truth kernel (the suspicion is the latter;
  Appendix~\ref{app:double_diag}); (b) whether the failure to
  converge in $\max|r-1|$ matters downstream. The latter has a
  clean answer: the PPG12 analysis applies fiducial cuts
  $|z_r|<60$\,cm and $|z_h|<60$\,cm, inside which the per-bin
  closure is already $1$--$2\%$. The stuck edge at
  $|z_r|\simeq 83$\,cm is well outside the analysis fiducial, so
  the iteration is good enough for the analysis use case.
\item \textbf{Implementation can be sped up}: each iteration
  currently rebuilds the MC reco histogram from scratch
  ($\mathcal{O}(N_\mathrm{events})$ per pass). Precomputing $z^2$
  and dropping the per-trial mask rebuild would give a
  $\sim\!10\times$ speedup with no change in result. Not a
  blocker; nice-to-have for the all-samples extension.
\item \textbf{DI-MC \texttt{vertexz\_truth\_mb} branch}: the
  \texttt{photon10\_double} and \texttt{jet12\_double}
  \texttt{bdt\_split.root} files currently fed to the efficiency
  pipeline predate the \texttt{vertexz\_truth\_mb} branch in the
  anatreemaker output, so $w(z_\mathrm{mb})$ is silently dropped
  (optional-pointer guard). The impact is expected sub-percent at
  1.5\,mrad and sub-percent at 0\,mrad; regenerating
  \texttt{bdt\_split.root} from the reprocessed
  \texttt{combined.root} (which does carry the branch) removes
  this caveat.
\end{itemize}


%% ==================================================================
\section{Cluster Size Under DI}
\label{sec:cluster_size}

Cluster size (EMCal tower count \nown plus the 7$\times$7
bounding-box widths \dietag, \diphig) is the basic shape property
that underlies every shower-shape-based discriminator used in
PPG12 (NPB, shower-shape BDT, width variables). This section
reports the DI-induced shifts, split by sample, selection stage,
and crossing-angle period.

\subsection{Sample-Level Reference (Pre-Cut)}
\label{subsec:cluster_di_ref}

Figure~\ref{fig:di_photon} shows the pure single vs.\ pure double
reference for photon10; Figure~\ref{fig:di_jet} for jet12. The
mixed curves at 7.9\% (1.5\,mrad) and 22.4\% (0\,mrad) are
bin-by-bin linear combinations of the two means, i.e.\
$\langle y\rangle_\mathrm{mix}(\ET) =
(1-f)\langle y\rangle_\mathrm{single}(\ET) +
 f \langle y\rangle_\mathrm{double}(\ET)$,
which avoids biasing the blend by the per-sample event counts.

\paragraph{Photons.}
DI shifts photon cluster size \emph{upwards} by
$\Delta\langle\nown\rangle \approx +0.15$ to $+0.22$ towers
above single across $\ET\in[15,30]$\,\GeV --- real but small.
Pile-up adds a few soft neighbours to an otherwise compact EM
shower. Bounding-box widths shift by $\lesssim 0.1$; the first-order
effect on photon shower width is at the sub-0.1-index level.
At 7.9\% (1.5\,mrad), the mixed curve is nearly indistinguishable
from single --- consistent with the shower-shape BDT being
pileup-robust in the 1.5\,mrad regime.

\begin{figure}[H]
\centering
\includegraphics[width=0.85\linewidth]{figures/cluster_size/cluster_size_di_photon10.pdf}
\caption{Photon DI reference: \texttt{photon10\_nom} (blue) vs.\
  \texttt{photon10\_double} (red), plus mixed at $f=7.9\%$ (green,
  1.5\,mrad) and $f=22.4\%$ (magenta, 0\,mrad). Pre-cut, both
  containers.}
\label{fig:di_photon}
\end{figure}

\paragraph{Jets.}
DI inflates \emph{jet} cluster size much more strongly:
$\Delta\langle\nown\rangle \approx +0.5$ to $+1.3$ towers in the
split container at $\ET>15$\,\GeV, roughly $4\times$ larger than the
photon response. Physically: a jet cluster already has multiple
soft neighbouring towers from accompanying hadrons, and DI adds
more of the same; the probability that the DI overlay lands on a
low-energy tower and pushes it over the 70\,MeV threshold is much
higher in a jet environment than for a clean photon. Because jets
enter both the fake-photon and anti-tight control regions of the
ABCD, this DI effect is a non-trivial input to the background-template
closure: in 0\,mrad (22.4\% DI) running, the jet template mean
cluster size shifts upward by $\sim\!0.3$--$0.4$ towers, which
affects the shower-shape BDT input distribution and therefore the
background rejection rate.

\begin{figure}[H]
\centering
\includegraphics[width=0.85\linewidth]{figures/cluster_size/cluster_size_di_jet12.pdf}
\caption{Jet DI reference: \texttt{jet12\_nom} vs.\
  \texttt{jet12\_double} in the same format as
  Fig.~\ref{fig:di_photon}. DI effect on jet cluster size is
  $\sim\!4\times$ larger than on photon, consistent with softer
  neighbouring towers getting pushed above the 70\,MeV threshold.}
\label{fig:di_jet}
\end{figure}

\subsection{Selection-Stage Progression (With Truth-Vertex Reweight)}
\label{subsec:cluster_selection}

With the nominal iterative truth-vertex reweight applied, three
selection stages are compared per period (baseline,
common-cut, tight-cut), re-binned to the
$\ET\in[19,21]$\,\GeV bin centre of the split container. The
numerical values at $\ET\simeq 20$\,\GeV are summarized in
Table~\ref{tab:cluster_numbers}.

\begin{table}[H]
\centering\small
\begin{tabular}{llrrrrrr}
\toprule
 & & \multicolumn{3}{c}{1.5\,mrad} & \multicolumn{3}{c}{0\,mrad}\\
\cmidrule(lr){3-5}\cmidrule(lr){6-8}
metric & sample & baseline & common & tight & baseline & common & tight \\
\midrule
$\langle\nown\rangle$   & data   & 15.86 & 11.01 & 9.92 & 14.37 & 11.77 & 9.98 \\
                         & photon & 10.44 & 10.39 & 10.21 & 10.56 & 10.51 & 10.32 \\
                         & jet    & 15.38 & 14.33 & 10.94 & 15.51 & 14.59 & 11.05 \\
\midrule
$\langle\dietag\rangle$  & data   &  5.65 &  3.90 & 3.63 &  5.10 &  4.15 & 3.66 \\
                         & photon &  3.78 &  3.77 & 3.72 &  3.83 &  3.82 & 3.76 \\
                         & jet    &  4.81 &  4.63 & 3.84 &  4.85 &  4.69 & 3.88 \\
\midrule
$\langle\diphig\rangle$  & data   &  4.24 &  4.03 & 3.82 &  4.19 &  4.06 & 3.80 \\
                         & photon &  3.93 &  3.91 & 3.88 &  3.93 &  3.91 & 3.88 \\
                         & jet    &  4.95 &  4.76 & 4.16 &  4.96 &  4.79 & 4.17 \\
\bottomrule
\end{tabular}
\caption{Cluster-size means in the bin centre
  $\ET\in[19,21]$\,\GeV, split container (\splitnode),
  truth-vertex-reweighted MC. Values read from the per-period
  TProfiles in
  \texttt{results/cluster\_size\_cuts\_PERIOD\_SAMPLE.root}
  (PERIOD = 1p5mrad or 0mrad). Photon pool is
  photon5+photon10+photon20; jet pool is
  jet5+jet8+jet12+jet20+jet30+jet40. DI blending is applied in the
  figures only, not in this numeric table.}
\label{tab:cluster_numbers}
\end{table}

Three observations:
\begin{itemize}
\item \textbf{Common cut removes the noise-like tower excess,
  preferentially from data.} At 1.5\,mrad and $\ET=20$\,\GeV,
  $\langle\nown\rangle$ drops from $15.86\to 11.01$ in data (a
  31\% decrease), $15.38\to 14.33$ in jet MC (only 7\%), and
  $10.44\to 10.39$ in photon MC (essentially unchanged). The
  common-cut's NPB+$w_\eta$-bound block targets
  noise/threshold contamination that is far more prevalent in data
  than in MC.
\item \textbf{Tight cut converges data, jet MC, and photon MC to
  the signal cluster size.} At $\ET=20$\,\GeV, 1.5\,mrad:
  data tight $9.92$, jet MC tight $10.94$, photon MC tight $10.21$
  --- a $\pm 5\%$ spread around a common $\sim\!10.4$-tower signal
  size. Photon-MC tight is the direct signal-MC comparator for
  data tight: the two agree to $\sim\!3\%$ on $\langle\nown\rangle$
  and $\langle\dietag\rangle$ across $\ET\in[10,30]$\,\GeV.
\item \textbf{DI mixing is sub-leading in the tight selection.}
  At 1.5\,mrad the DI-blend shift in $\langle\nown\rangle$ is below
  0.1 tower for every selection. At 0\,mrad (22.4\% DI) the jet-MC
  baseline shifts upward by $\sim\!0.3$ towers, but drops to below
  0.1 tower after the tight cut. Photon-MC DI effect is below
  0.1 tower at every stage and every period.
\end{itemize}

\begin{figure}[H]
\centering
\begin{subfigure}[t]{0.48\textwidth}
\includegraphics[width=\textwidth]{figures/cluster_size/cluster_size_period_baseline_split_1p5mrad.pdf}
\caption{1.5\,mrad baseline}
\end{subfigure}\hfill
\begin{subfigure}[t]{0.48\textwidth}
\includegraphics[width=\textwidth]{figures/cluster_size/cluster_size_period_baseline_split_0mrad.pdf}
\caption{0\,mrad baseline}
\end{subfigure}
\caption{Baseline cluster-size per period, split container
  (\splitnode), data vs.\ pooled photon MC vs.\ pooled jet MC,
  with truth-vertex-reweighted MC. The data-jet-MC gap is larger
  at 1.5\,mrad ($+0.5$ towers) than at 0\,mrad ($-1.1$ towers);
  the period difference at baseline largely absorbs into the tight
  selection.}
\end{figure}

\begin{figure}[H]
\centering
\begin{subfigure}[t]{0.48\textwidth}
\includegraphics[width=\textwidth]{figures/cluster_size/cluster_size_period_tight_nosplit_1p5mrad.pdf}
\caption{1.5\,mrad tight (no-split)}
\end{subfigure}\hfill
\begin{subfigure}[t]{0.48\textwidth}
\includegraphics[width=\textwidth]{figures/cluster_size/cluster_size_period_tight_nosplit_0mrad.pdf}
\caption{0\,mrad tight (no-split)}
\end{subfigure}
\caption{Tight-selection cluster-size per period in the no-split
  container (\nosplitnode). After the tight cut, data, jet MC, and
  photon MC converge to the photon signal size within $\pm 5\%$
  across both periods; DI effects on tight jet MC are at the
  sub-tower level.}
\end{figure}

The DI effect on the ABCD background template is therefore most
relevant at the \emph{common-cut} stage, where the 22.4\% jet-MC
cluster-size shift feeds the shower-shape BDT distribution, and
at 0\,mrad more than at 1.5\,mrad. After the full tight
selection, the DI effect on the background template is absorbed
to below the $\sim\!1\%$ level on $\langle\nown\rangle$, which is
well within the systematic budget of the background-shape cross-check.


%% ==================================================================
\section{Energy Response Under DI}
\label{sec:energy_response}

The photon energy response $R = \ET^\mathrm{reco}/\pT^\mathrm{truth}$
is measured bin-by-bin in truth \pT across the single, double, and
mixed MC, both with and without the truth-vertex reweight. The
extraction is driven by \texttt{efficiencytool/compare\_energy\_response.C},
which builds six sample sets per \pT bin
(single, double, mixed\_0mrad, mixed\_0mrad\_vtxrw, mixed\_1p5mrad,
mixed\_1p5mrad\_vtxrw) and fits each to a double-sided Crystal-Ball
(DSCB) function to extract the peak $\mu$ and resolution $\sigma$.
The blended vertex scan used for the MC vertex reweight is built
internally from the \texttt{photon10\_nom} and \texttt{photon10\_double}
vertex distributions.

\begin{figure}[H]
\centering
\begin{subfigure}[t]{0.48\textwidth}
\includegraphics[width=\textwidth]{../plotting/figures/energy_response/resp_mean_common_cut_pileup.pdf}
\caption{Response mean $\mu$, DI comparison}
\end{subfigure}\hfill
\begin{subfigure}[t]{0.48\textwidth}
\includegraphics[width=\textwidth]{../plotting/figures/energy_response/resp_sigma_common_cut_pileup.pdf}
\caption{Response width $\sigma$, DI comparison}
\end{subfigure}
\caption{Photon energy response $R=\ET^\mathrm{reco}/\pT^\mathrm{truth}$
  at the common-cut level, across single, double, and mixed
  (0\,mrad / 1.5\,mrad) MC without the truth-vertex reweight.
  Left: DSCB peak $\mu(\pT)$; right: DSCB width $\sigma(\pT)$.
  Double MC shifts the peak by $\lesssim 0.5\%$ and broadens the
  width by $\lesssim 10\%$ at $\pT\sim 20$\,\GeV; the mixed curves
  sit between single and double at the appropriate fractions.}
\label{fig:resp_pileup}
\end{figure}

\begin{figure}[H]
\centering
\begin{subfigure}[t]{0.48\textwidth}
\includegraphics[width=\textwidth]{../plotting/figures/energy_response/resp_mean_common_cut_vtxrw.pdf}
\caption{Response mean $\mu$, with truth-vertex reweight}
\end{subfigure}\hfill
\begin{subfigure}[t]{0.48\textwidth}
\includegraphics[width=\textwidth]{../plotting/figures/energy_response/resp_sigma_common_cut_vtxrw.pdf}
\caption{Response width $\sigma$, with truth-vertex reweight}
\end{subfigure}
\caption{Photon energy response with the iterative truth-vertex
  reweight applied on the mixed samples (mixed\_0mrad\_vtxrw,
  mixed\_1p5mrad\_vtxrw). The truth-vertex reweight further
  reduces the residual DI-induced mean shift to below $0.3\%$
  across the full \pT range. Layout as in Fig.~\ref{fig:resp_pileup}.}
\label{fig:resp_vtxrw}
\end{figure}

\begin{figure}[H]
\centering
\begin{subfigure}[t]{0.48\textwidth}
\includegraphics[width=\textwidth]{../plotting/figures/energy_response/resp_fit_qc_grid_dscb_single_tight_iso.pdf}
\caption{Single MC, tight+iso}
\end{subfigure}\hfill
\begin{subfigure}[t]{0.48\textwidth}
\includegraphics[width=\textwidth]{../plotting/figures/energy_response/resp_fit_qc_grid_dscb_double_tight_iso.pdf}
\caption{Double MC, tight+iso}
\end{subfigure}\\[6pt]
\begin{subfigure}[t]{0.48\textwidth}
\includegraphics[width=\textwidth]{../plotting/figures/energy_response/resp_fit_qc_grid_dscb_mixed_0mrad_tight_iso.pdf}
\caption{Mixed 0\,mrad, tight+iso}
\end{subfigure}\hfill
\begin{subfigure}[t]{0.48\textwidth}
\includegraphics[width=\textwidth]{../plotting/figures/energy_response/resp_fit_qc_grid_dscb_mixed_1p5mrad_tight_iso.pdf}
\caption{Mixed 1.5\,mrad, tight+iso}
\end{subfigure}
\caption{Per-bin DSCB fits to the photon energy response at the
  tight+iso selection level for single, double, and mixed MC.
  The DSCB parameters are stable across bins; the single-vs-double
  spread is dominated by the $\mu$ and $\sigma$ shifts of
  Figs.~\ref{fig:resp_pileup}--\ref{fig:resp_vtxrw}.}
\label{fig:resp_dscb_grid}
\end{figure}

The DI-induced peak shift of $\lesssim 0.5\%$ and width
broadening of $\lesssim 10\%$ at the common-cut level sits well
within the existing energy-scale systematic budget of PPG12 and is
further reduced after the iterative truth-vertex reweight. The
energy response is therefore not a dominant DI systematic for the
cross-section measurement; the leading DI sensitivity remains on
the efficiency stage (\S\ref{sec:eff_impact}) and the ABCD
background template (\S\ref{sec:cluster_size}).


%% ==================================================================
\section{Eta-Edge Migration Under Pileup}
\label{sec:eta_edge}

The eta-edge migration study ($|\eta|<0.7$ fiducial) exploits the
fact that DI shifts the cluster $\eta$ via
Eq.~\eqref{eq:deta_vs_dz}, so a cluster whose truth $\eta$ is
just inside the fiducial edge can migrate outside the fiducial on
reconstruction (producing a signal loss) and, symmetrically, a
cluster whose truth $\eta$ is just outside can migrate inside
(producing a background fake). In the nominal 1.5\,mrad period the
effect is sub-leading: the vertex-shift RMS is small
($\sigma_{\Delta z}\simeq 7.6$\,cm single,
$\sigma_{\Delta z}\simeq 34.5$\,cm double; only 7.9\% of clusters
draw from the wider double distribution), which translates to
$\sigma_{\Delta\eta}\simeq 0.08$ in the worst case. The fiducial
edge at $|\eta|=0.7$ is $\gtrsim 5\sigma_{\Delta\eta}$ from the
barrel edge at $|\eta|\simeq 1$, so the out-migration rate is
percent-level. At 0\,mrad the effect is proportionally larger
($\sigma_{\Delta\eta}\simeq 0.37$ averaged), but the 0\,mrad
period is not part of the nominal published cross-section.

The full eta-edge migration analysis (fake rates, signal leakage,
ABCD $R$-factor bias from edge effects, net sideband flow) is
documented in \texttt{wiki/concepts/eta-edge-migration.md} and the
companion report in \texttt{reports/figures/} (figures
\texttt{eta\_migration\_abcd\_fakes\_double.pdf},
\texttt{eta\_migration\_matrix\_double.pdf} and peers). For the
purposes of this DI note, the relevant observation is that the
fiducial-migration effect is absorbed by the response-matrix
unfolding used by PPG12 and does not need a separate DI correction.
It is flagged here only because the underlying mechanism
(vertex-shift-induced $\Delta\eta$) is shared with the other DI
effects analysed in this note.


%% ==================================================================
\section{Summary, Systematics, and Recommendations}
\label{sec:summary}

\subsection{Nominal Recipe}

The PPG12 DI programme converges on the following recipe for the
nominal cross-section extraction:

\begin{enumerate}
\item \textbf{Efficiency: single-MC only, with track-ID-only
  matching (no $\DR<0.1$ cut).} The apparent double-MC efficiency
  drop is dominated by a matching artifact
  (\S\ref{subsec:decomposition}); removing $\DR<0.1$ eliminates
  the artifact and reveals the genuine pileup degradation. Under
  track-ID matching the total photon-ID efficiency drops from
  57.4\% (single MC) to 52.0\% (mixed 0\,mrad), a
  $1-52.0/57.4 = 9.4\%$ relative reduction concentrated in the
  reco and BDT stages. At the PPG12 1.5\,mrad working point
  ($f_\mathrm{DI}=7.9\%$), the integrated DI effect on the
  cross section is $1.6\%$ (reco) $+ 1.0\%$ (BDT) $+ 0.2\%$ (iso)
  $= 2.8\%$ pre-correction (see \S\ref{subsec:eff_implications}).
\item \textbf{Vertex reweighting: iterative truth-vertex
  $w(z)$ applied per truth vertex.} The production reco-level
  $f(z_r)$ structurally cannot close the individual truth
  vertices of DI events; the iterative method
  (\S\ref{sec:truth_vertex}) closes the truth fiducial to
  $1$--$2\%$ inside $|z_h|<60$\,cm at 1.5\,mrad and converges
  cleanly at 0\,mrad. It removes the $+1.4\%$ cross-section bias
  at 1.5\,mrad that the reco-level method alone could not address.
  The reco-level reweight remains available as a
  mutually-exclusive fallback; the truth-vertex variant is the
  nominal going forward.
\item \textbf{Shower-shape templates: mixed MC for cross-check
  only, not as the nominal template.} The data-vs-mixed-MC
  validation (\S\ref{sec:shower_shape}) confirms that the mixed
  MC improves data-MC \chidf by up to $20.8\times$ for
  pileup-sensitive observables at high \ET; the improvement does
  not hold at low \ET (pt0, $10$--$14$\,\GeV) and the mixed MC
  degrades there. The analysis uses single MC for the nominal BDT
  evaluation and quotes the nominal-vs-mixed spread as a
  shower-shape systematic.
\item \textbf{ABCD background template: photon-MC-free, built
  from the anti-tight region in data.} The DI-induced jet-MC
  cluster-size shift at 0\,mrad is sub-leading after the tight
  selection (\S\ref{sec:cluster_size}); at 1.5\,mrad it is
  negligible.
\end{enumerate}

\subsection{Residual Systematic}

With the above recipe in place, the residual DI-related
systematic at 1.5\,mrad has the following ingredients:

\begin{itemize}
\item \textbf{Truth-vertex reweight residual}: $1$--$2\%$ per bin
  inside $|z_h|<60$\,cm, sub-percent integrated over the analysis
  fiducial.
\item \textbf{Efficiency pileup effect}: 2.8\% integrated at
  1.5\,mrad (reco 1.6\% + BDT 1.0\% + iso 0.2\%; see
  \S\ref{subsec:eff_implications}); an explicit
  single-vs-mixed-MC efficiency systematic at the cluster-weighted
  pileup fraction covers this contribution.
\item \textbf{Shower-shape mismodelling at low \ET}: picked up by
  the nominal-vs-mixed-MC spread on the BDT templates, explicitly
  quantified in the shower-shape systematic.
\item \textbf{ABCD background-template DI shift}: sub-percent at
  1.5\,mrad after tight selection; absorbed by the
  background-template systematic.
\item \textbf{Missing \texttt{vertexz\_truth\_mb} branch on current
  DI \texttt{bdt\_split.root}}: sub-percent drop of the
  secondary-vertex factor $w(z_\mathrm{mb})$;
  expected sub-percent impact at 1.5\,mrad and sub-percent at
  0\,mrad. Removed by regenerating the DI \texttt{bdt\_split.root}
  files.
\end{itemize}

The dominant pre-existing cross-section bias at 1.5\,mrad,
$+1.4\%$ from the truth-fiducial mistreatment of DI events under
$f(z_r)$, is eliminated by the truth-vertex reweight.

\subsection{Cross-Section Impact Summary}

\begin{table}[H]
\centering
\caption{DI-related effects on the 1.5\,mrad cross section and
  their treatment. Percentages are integrated over the analysis
  fiducial unless otherwise noted.}
\label{tab:xsec_summary}
\small
\begin{tabular}{p{4.5cm}p{3.5cm}p{3cm}p{3cm}}
\toprule
Effect & Pre-fix bias (1.5\,mrad) & After nominal recipe & Treatment \\
\midrule
Truth fiducial mistreatment under $f(z_r)$ alone & $+1.4\%$ & $\lesssim 0.1\%$ & Iterative truth-vertex reweight \\
Efficiency matching artifact ($\DR<0.1$) & $-3.5\%$ baseline in single MC, dominant in double MC &  absorbed into nominal single-MC efficiency & Track-ID-only matching \\
Efficiency pileup (reco + BDT + iso combined) & $2.8\%$ (1.6+1.0+0.2) & covered by single-vs-mixed-MC systematic & Single MC nominal; mixed MC systematic \\
ABCD jet-template DI shift & few \% on background shape at common-cut stage & sub-\% on tight-selection yield & Absorbed by background-template systematic \\
Shower-shape low-\ET mismodelling & several $\sigma$ on \chidf & covered by mixed-vs-nominal systematic & Nominal-vs-mixed BDT systematic \\
Missing $w(z_\mathrm{mb})$ on DI trees & --- & sub-\% & Regenerate DI \texttt{bdt\_split.root} \\
\bottomrule
\end{tabular}
\end{table}

\subsection{Open Items}

\begin{itemize}
\item \textbf{jet12 extension of the truth-vertex fit.} Current
  iterative reweight uses photon10-only. Extending to the full
  4-sample blend (photon10/photon10\_double + jet12/jet12\_double)
  is straightforward but the closure numbers will move and should
  be remeasured.
\item \textbf{DI \texttt{bdt\_split.root} regeneration.} The
  current \texttt{photon10\_double} and \texttt{jet12\_double}
  \texttt{bdt\_split.root} files predate the
  \texttt{vertexz\_truth\_mb} branch; regenerating from the
  reprocessed \texttt{combined.root} removes the silent
  secondary-vertex drop.
\item \textbf{NPB purity 0 vs.\ 1.5\,mrad split.} Deferred; the
  NPB-cut systematic is currently evaluated integrated over both
  periods. A per-period split would sharpen the pileup sensitivity.
\item \textbf{Iterative-reweight speed optimization.} $\sim\!10\times$
  speedup available via pre-computed $z^2$ and dropping the
  per-trial mask rebuild; nice-to-have for multi-sample extensions.
\item \textbf{Pipeline consolidation.} 14-way duplication of the
  run28 condor submit scripts + the hardcoded
  \texttt{total\_line=10000} cap are both tracked in
  \texttt{condor\_pipeline\_review\_run28.md}; neither blocks the
  physics at current precision.
\end{itemize}


%% ==================================================================
\appendix
\section{DI Reco--Truth Kernel}
\label{app:double_diag}

The geometric lower bound on truth-only reweights (referenced in
Appendix~\ref{app:discarded}) can be read off the widths directly
for the jet12\_double sample at 1.5\,mrad:

\begin{itemize}\setlength\itemsep{1pt}
\item 1.5\,mrad data reco width: $\sigma \simeq 22.3$\,cm
\item jet12 single-interaction MC truth width ($z_h$): $\sigma \simeq 57.6$\,cm
\item jet12\_double MC truth width ($z_h$, primary): $\sigma \simeq 62.7$\,cm
\item jet12\_double MC reco width ($z_r$):
      $\sigma \simeq 48.6 \approx 62.7/\sqrt 2$ (averaging)
\item Conditional kernel $P(z_r|z_h=0)$ RMS (with $z_\mathrm{mb}$
      marginalized): $\sim\!34$\,cm
\end{itemize}

Figure~\ref{fig:double_diag} shows the kernel slice directly,
compared to the 1.5\,mrad data target. Even if the primary truth
vertex is forced to $z_h=0$, the reco distribution from a
double-interaction event cannot be narrower than the
$\sim\!34$\,cm kernel width --- larger than the 22\,cm target.
This is the structural reason the truth-only methods of
Appendix~\ref{app:discarded} fail on DI MC, and the structural
reason the iterative method of \S\ref{sec:truth_vertex} needs the
$w(z_h)\,w(z_\mathrm{mb})$ factorization rather than a $w$-on-$z_r$
construction.

\begin{figure}[H]
\centering
\includegraphics[width=\linewidth]{figures/vertex_reweight_alt_double_diag.pdf}
\caption{Conditional reco kernel $P(z_r|z_h)$ in jet12\_double MC
  at 1.5\,mrad (with $z_\mathrm{mb}$ marginalized), with the
  1.5\,mrad data reco $z_r$ shape overlaid. The kernel's
  $\sim\!34$\,cm RMS exceeds the $\sim\!22$\,cm data peak, which
  is the geometric lower bound on truth-only reweighting.}
\label{fig:double_diag}
\end{figure}


\section{Discarded Reweighting Methods}
\label{app:discarded}

Before the iterative method of \S\ref{sec:truth_vertex} was
settled on, three truth-only reweight constructions and one
closed-form 2D construction were explored. All four were
discarded; the essentials:

\paragraph{$w_1$: direct ratio.}
$w_1(z_h) = D(z_r)/T(z_h)$ with the common binning identifying
$z_r$ bins with $z_h$ bins. Breaks down on the DI sample because
$P(z_r|z_h)$ is wide: the LHS has a narrow data target while the
RHS integrates over a wide kernel.

\paragraph{$w_2$: conditional projection.}
$w_2(z_h) = \mathbb{E}[f(z_r)\,|\,z_h]$, the expectation of the
production reco weight conditional on the primary truth vertex.
Also bounded from below by the DI kernel width
(Appendix~\ref{app:double_diag}).

\paragraph{$w_\mathrm{iter,DA}$: D'Agostini iterative unfolding.}
Iterative Bayesian unfolding using the MC response matrix
$R_{ij}=P(z_r^i|z_h^j)$. Hits the same geometric floor at
1.5\,mrad --- the DI reco-truth kernel is too wide to invert
within the data-domain support.

\paragraph{Closed-form 2D Gaussian reweight.}
$w(z_h,z_\mathrm{mb}) \propto \exp[-\tfrac{1}{2}(z_h^2 + z_\mathrm{mb}^2)\,\Delta]$
with $\Delta = 1/\sigma_\mathrm{lum}^2 - 1/\sigma_\mathrm{MC}^2$.
This acts on both the symmetric and antisymmetric coordinates
simultaneously and closes all three vertex axes; however, it
assumes a single Gaussian shape for the data luminous region. At
1.5\,mrad the data has a $\pm$-asymmetric mid-tail bump that a
symmetric Gaussian cannot reproduce, so the closed form was
superseded by the iterative method that discovers $w(z)$ from the
data directly without any Gaussian assumption.

\begin{figure}[H]
\centering
\begin{subfigure}[t]{0.48\textwidth}
\includegraphics[width=\textwidth]{figures/vertex_reweight_alt_weights.pdf}
\caption{Recovered $w(z)$: discarded methods}
\end{subfigure}\hfill
\begin{subfigure}[t]{0.48\textwidth}
\includegraphics[width=\textwidth]{figures/vertex_reweight_alt_shapes.pdf}
\caption{Truth-vertex shapes: discarded methods}
\end{subfigure}\\[6pt]
\begin{subfigure}[t]{0.48\textwidth}
\includegraphics[width=\textwidth]{figures/vertex_reweight_alt_closure.pdf}
\caption{Closure: discarded methods}
\end{subfigure}\hfill
\begin{subfigure}[t]{0.48\textwidth}
\includegraphics[width=\textwidth]{figures/vertex_reweight_alt_jet_vs_photon_closure.pdf}
\caption{Jet vs.\ photon closure: discarded methods}
\end{subfigure}
\caption{Diagnostic plots for the discarded truth-only reweight
  methods. All four methods hit a geometric floor on the DI sample
  at 1.5\,mrad set by the reco-truth kernel width.}
\end{figure}


\section{TruthVertexRecoCheck Baseline}
\label{app:nonull}

Before building any truth-based reweight, a sanity check on the
effect of dropping MC events whose reco MBD vertex is the $-9999$
sentinel was performed. The concern: an MC event that loses its
reco vertex may be physically unusual (e.g.\ with a shifted truth
vertex $z_h$), so dropping it could skew the efficiency in a
$z_h$-dependent way.

\begin{table}[H]
\centering
\begin{tabular}{lrrrr}
\toprule
Sample & $N_\mathrm{total}$ & $N_\mathrm{reco}$ & $N_\mathrm{no\,reco}$ & fraction no reco \\
\midrule
photon10           & 9{,}998{,}562 & 7{,}716{,}585 & 2{,}281{,}977 & \textbf{22.82\%} \\
photon10\_double   & 9{,}997{,}075 & 9{,}491{,}223 &   505{,}852   & \textbf{5.06\%} \\
jet12              & 10{,}000{,}000 & 7{,}689{,}530 & 2{,}310{,}470 & \textbf{23.10\%} \\
jet12\_double      & 9{,}997{,}000 & 9{,}485{,}384 &   511{,}616   & \textbf{5.12\%} \\
\bottomrule
\end{tabular}
\caption{Fraction of events without an MBD reco vertex. The
  double-interaction productions have $\sim\!4.5\times$ higher
  MBD vertex-finding efficiency than their single counterparts
  (two concurrent collisions give the MBD twice the chance of
  producing a vertex).}
\label{tab:nonull_counts}
\end{table}

\begin{figure}[H]
\centering
\includegraphics[width=0.95\linewidth]{figures/truth_vertex_reco_check.pdf}
\caption{Primary truth vertex $z_h$ distribution for events with
  (blue) and without (red dashed) a reconstructed MBD vertex
  $z_r$, area-normalized for shape comparison. Columns: single /
  double interaction. Rows: photon10 / jet12. Dotted vertical
  lines: $|z_h|=30$\,cm reference.}
\label{fig:nonull_shape}
\end{figure}

\begin{figure}[H]
\centering
\includegraphics[width=0.95\linewidth]{figures/truth_vertex_reco_check_ratio.pdf}
\caption{Ratio (no-reco-vertex shape)/(reco-vertex shape) vs.\
  $z_h$, zoomed to $[0.75,1.25]$. All samples show a $\sim\!\pm 5\%$
  north/south asymmetry inside $|z_h|<30$\,cm, consistent with a
  small MBD geometric bias.}
\label{fig:nonull_ratio}
\end{figure}

Inside $|z_h|<30$\,cm the two shapes agree to within $\pm 5\%$ with
a small $z_h$-asymmetry that is the same in all four samples and
therefore a detector-level effect (MBD vertex finder slightly more
efficient at positive $z$). Dropping the no-reco-vertex events is
\emph{approximately} harmless --- both the efficiency numerator and
denominator are scaled by the same factor, and the residual $z$
dependence of the drop rate is small. However, dropping $\sim\!22\%$
of the single-interaction MC sample is still a lot of statistics,
and it makes the efficiency measurement implicitly conditional on a
valid reco vertex being found. That was the original motivation for
exploring a truth-level reweight in the first place.


\section{File and Code Inventory}
\label{app:code}

\paragraph{Code.}
\begin{itemize}\setlength\itemsep{1pt}
\item \texttt{anatreemaker/source/CaloAna24.\{h,cc\}} --- writes the
  \texttt{vertexz\_truth} ($z_h$) and
  \texttt{vertexz\_truth\_second} ($z_\mathrm{mb}$, read in C++ as
  \texttt{vertexz\_truth\_mb}) branches on the slimtree.
\item \texttt{efficiencytool/ShowerShapeCheck.C} --- shower-shape
  analysis with the \texttt{mix\_weight} +
  \texttt{vtxscan\_sim\_override} DI-blend options.
\item \texttt{efficiencytool/run\_showershape\_double.sh} ---
  two-pass full-GEANT blending pipeline.
\item \texttt{efficiencytool/RecoEffCalculator\_TTreeReader.C} ---
  efficiency calculator (\texttt{trkID}-only matching;
  \texttt{SetUseWeightedEvents} fixed).
\item \texttt{efficiencytool/compare\_energy\_response.C} ---
  energy response $\mu(\pT)$, $\sigma(\pT)$ under DSCB across
  single/double/mixed MC with optional vertex reweight.
\item \texttt{efficiencytool/truth\_vertex\_reweight/} --- standalone
  Python tool for the iterative method:
  \begin{itemize}\setlength\itemsep{0pt}
  \item \texttt{fit\_core.py} --- cached sample loader,
    \texttt{mixed\_mc\_histogram} forward model, Poisson $\chi^2$.
  \item \texttt{fit\_truth\_vertex\_reweight.py} --- driver
    producing \texttt{reweight.root} and \texttt{fit\_result.yaml}.
  \item \texttt{plot\_closure.py}, \texttt{plot\_comparisons.py} ---
    closure and cross-period diagnostic plots.
  \item \texttt{config.yaml} --- sample paths, run ranges, fit
    parameters.
  \end{itemize}
\item \texttt{efficiencytool/TruthVertexRecoCheck.C} ---
  Appendix~\ref{app:nonull} supporting macro.
\item \texttt{efficiencytool/calc\_pileup\_range.C} --- MBD-based
  pileup-fraction calculator.
\item \texttt{efficiencytool/MbdPileupHelper.h} --- MBD pileup
  metric helper (avg / prod / max $\sigma_t$ etc.).
\item \texttt{efficiencytool/DoubleInteractionCheck.C} --- toy
  vertex-shift simulation (not part of the nominal DI treatment;
  see \texttt{.claude/rules/double-interaction.md}).
\end{itemize}

\paragraph{Config fields.}
\begin{itemize}\setlength\itemsep{1pt}
\item \texttt{truth\_vertex\_reweight\_on} (0/1): enables the
  iterative truth-vertex reweight.
\item \texttt{truth\_vertex\_reweight\_file}: path to
  \texttt{reweight.root} containing \texttt{h\_w\_iterative}.
\item \texttt{vertex\_reweight\_on} (0/1): legacy reco-level
  $f(z_r)$; mutually exclusive with \texttt{truth\_vertex\_reweight\_on}.
\item \texttt{vertex\_scan\_data\_file}: unrestricted data vertex
  distribution for the toy sim (DI cross-check only).
\end{itemize}

\paragraph{Output files.}
\begin{itemize}\setlength\itemsep{1pt}
\item \texttt{efficiencytool/truth\_vertex\_reweight/output/\{1p5mrad,0mrad\}/reweight.root}
  --- per-period $w(z)$ and diagnostic histograms.
\item \texttt{efficiencytool/truth\_vertex\_reweight/output/\{1p5mrad,0mrad\}/fit\_result.yaml}
  --- fit metadata and convergence statistics.
\item \texttt{results/MC\_efficiencyshower\_shape\_\{sample\}\_\{suffix\}\_vtxscan.root}
  --- Pass-1 blended vertex distribution.
\item \texttt{results/MC\_efficiencyshower\_shape\_\{sample\}\_combined[\_inclusive]\_\{suffix\}.root}
  --- Pass-2 merged full-analysis output.
\end{itemize}

\paragraph{Known source-plot cosmetic issues (reused as-is).}
Two classes of known source-PDF cosmetic issues are reused from the
upstream plotting code without post-editing (this note preserves the
existing figures verbatim rather than regenerating them):
\begin{itemize}\setlength\itemsep{1pt}
\item The per-bin DSCB response grids in \S\ref{sec:energy_response}
  (Fig.~\ref{fig:resp_dscb_grid}) can show slight legend/data overlap
  in the top-left panel on some pages; this is a \texttt{matplotlib}
  auto-layout artifact in the upstream \texttt{compare\_energy\_response.C}
  grid plotter and does not affect the numerical fits.
\item The truth-vertex iterative-closure and method-comparison
  figures (Figs.~\ref{fig:iter_two_period}--\ref{fig:sigma_scan})
  carry Python \texttt{fig.suptitle()} banners inherited from
  \texttt{truth\_vertex\_reweight/plot\_closure.py}. These banners
  duplicate information already captured in the LaTeX captions but
  are kept unedited to preserve byte-level reproducibility against
  the \texttt{truth\_vertex\_reco\_check.pdf} source report.
\end{itemize}
These are flagged here for transparency and do not reflect errors in
the physics content.

\end{document}
